\documentclass[11pt]{article}

\usepackage[a4paper,margin=1in]{geometry}
\usepackage{amsmath,amssymb,amsthm,mathtools,bm}
\usepackage{booktabs,array,multirow,tabularx}
\usepackage{graphicx}
\usepackage{xcolor}
\usepackage{microtype}
\usepackage{enumitem}
\usepackage{url}
\usepackage[colorlinks=true,linkcolor=blue!55!black,
  citecolor=blue!55!black,urlcolor=blue!55!black]{hyperref}
\usepackage[nameinlink,noabbrev]{cleveref}

\theoremstyle{plain}
\newtheorem{theorem}{Theorem}[section]
\newtheorem{proposition}[theorem]{Proposition}
\newtheorem{lemma}[theorem]{Lemma}
\newtheorem{corollary}[theorem]{Corollary}
\theoremstyle{definition}
\newtheorem{definition}[theorem]{Definition}
\newtheorem{assumption}[theorem]{Assumption}
\newtheorem{ccopalgorithm}{Algorithm}
\theoremstyle{remark}
\newtheorem{remark}[theorem]{Remark}

\newcommand{\R}{\mathbb{R}}
\newcommand{\B}{\mathcal{B}}
\newcommand{\M}{\mathcal{M}}
\newcommand{\V}{\mathcal{V}}
\newcommand{\Q}{\mathcal{Q}}
\newcommand{\K}{\mathcal{K}}
\newcommand{\D}{\mathcal{D}}
\newcommand{\G}{\mathcal{G}}
\newcommand{\Grad}{\operatorname{Grad}}
\newcommand{\Div}{\operatorname{Div}}
\newcommand{\tr}{\operatorname{tr}}
\newcommand{\cof}{\operatorname{cof}}
\newcommand{\sym}{\operatorname{sym}}
\newcommand{\range}{\operatorname{range}}
\newcommand{\kernel}{\operatorname{ker}}
\newcommand{\diag}{\operatorname{diag}}
\newcommand{\rms}{\operatorname{RMS}}
\newcommand{\dt}{\Delta t}
\newcommand{\Th}{\Theta}
\newcommand{\half}{\tfrac12}
\newcommand{\norm}[2][]{\left\lVert #2\right\rVert_{#1}}
\newcommand{\ip}[3][]{\left\langle #2,#3\right\rangle_{#1}}
\newcommand{\figroot}{results}
\newcommand{\safegraphic}[2][0.94\linewidth]{%
  \IfFileExists{#2}{\includegraphics[width=#1]{#2}}{%
    \fbox{\parbox[c][35mm][c]{0.90\linewidth}{\centering
    Figure file not found:\\[-1mm]\scriptsize\path{#2}}}}}

\setlength{\emergencystretch}{3em}

\title{Configuration Consistency in Lagrangian Pressure Projection: Terminal Incompressibility and Pressure Work}
\author{Authors withheld for anonymous review}
\date{}

\begin{document}
\maketitle

\begin{abstract}
Standard semi-implicit Lagrangian projection methods evaluate incompressibility and pressure action on one geometric stage. We show that this construction hides a configuration-assignment problem. The accepted velocity constraint belongs to the terminal configuration, whereas the pressure action belongs to a separately selected material configuration. We formulate this separation as configuration-consistent oblique projection (CCOP), with cross response \(D_TG_\Theta\) and one scalar pressure equation. The parameter \(\Theta\) assigns the pressure operator geometrically while preserving the terminal incompressibility target. Every admissible placement satisfies an exact discrete pressure work--energy identity with its represented action, while placement determines the value and physical accuracy of that work. We establish the projection structure, relative-deformation solvability, and a local branch of the coupled step. On a fixed spatial discretization, midpoint action placement cancels the leading configuration-induced velocity defect for the force-free step and the second-order conservative-force predictor used here. Affine tests isolate terminal closure, projector obliqueness, and pressure work; a matched ISPH comparison shows their long-time consequences. A strictly nonaffine differential vortex verifies transported operators and complete evolution under spatially varying material shear. CCOP therefore makes the finite-step relation between terminal closure and pressure work explicit with one pressure stage.
\end{abstract}

\noindent\textbf{Keywords.}
Lagrangian incompressible flow, pressure projection, constrained mechanics,
Piola transform, generalized moving least squares, oblique projection.

\noindent\textbf{AMS subject classifications.}
65M12, 65M25, 65M75.

\section{Introduction}
\label{sec:introduction}

Pressure projection is one of the most effective ways to impose incompressibility: a predicted velocity is corrected through a scalar pressure solve, and the pressure equation can be embedded in finite-difference, finite-element, particle, or meshless discretizations \cite{Chorin1968,Temam1969,GuermondMinevShen2006}. In a Lagrangian method, this structure is especially attractive because the particles themselves carry the evolving material geometry. This moving geometry, however, introduces two distinct geometric roles that coincide on a fixed domain. The incompressibility constraint is associated with the accepted terminal configuration, whereas the pressure action is represented on a material configuration through a configuration-dependent Piola map. These two choices govern different properties: the former determines the represented incompressibility condition satisfied by the accepted state, while the latter controls the pressure impulse and its kinetic work. A Lagrangian projection must therefore specify not only which constraint is closed, but also where the multiplier acts.

Particle projection methods have advanced along several complementary directions, including pressure-source construction, boundary treatment, spatial consistency, and time integration \cite{CumminsRudman1999,XuStansbyLaurence2009,KhayyerGotoh2011,GotohKhayyerGotoh2024}. Nair and Tomar~\cite{NairTomar2015} connected incompressibility to the deformation and volume state of the updated particle configuration through a deformation-gradient-based pressure formulation. Consistent high-order ISPH improves the spatial pressure operators and boundary treatment together \cite{ZhangEtAl2024}, while Matsunaga~\cite{Matsunaga2025} develops trajectory-based pressure equations with high-order Runge--Kutta integration.  These developments address complementary levels of projection design: volume-state representation, spatial consistency, and temporal integration. Here we address another structural level---the configuration assignment underlying the projection itself. At the continuum level we require incompressibility on the accepted terminal configuration, so that the corrected velocity is tangent to the terminal volume-preserving state; the particle scheme closes its explicitly identified pressure-test representation of that tangent condition. The remaining question is where the pressure gradient should act along the material path. This leads naturally to a projection with terminal divergence and an independently placed intermediate pressure action.

A useful structural analogue appears in constrained mechanics. Discrete mechanics and variational integrators assign forces and constraints distinct roles in an update \cite{MarsdenWest2001}. RATTLE and discrete null-space methods organize the same distinction within structure-preserving time integration \cite{Andersen1983,Betsch2005}. Related constructions also appear in incompressible particle and geometric fluid discretizations \cite{SuzukiKoshizukaOka2007,PavlovEtAl2011,GawlikEtAl2011}. CCOP develops the corresponding configuration assignment inside a pressure-projection/PPE factorization: terminal geometry carries the accepted constraint, and a selected material geometry carries the pressure multiplier action.

Let $D_T$ denote the tangent of the accepted terminal incompressibility
constraint and let $G_\Theta$ denote the Piola pressure action represented on a
real material action configuration.  The frozen pressure correction has the
mixed form
\begin{equation}
 \begin{bmatrix}
   I & \Delta t\,G_\Theta\\
   D_T & 0
 \end{bmatrix}
 \begin{bmatrix}u_T\\ p\end{bmatrix}
 =
 \begin{bmatrix}u^\star\\0\end{bmatrix}.
 \label{eq:intro_saddle_factorization}
\end{equation}
Schur elimination gives
\begin{equation}
 D_TG_\Theta p=\Delta t^{-1}D_Tu^\star,
 \label{eq:intro_cross_response}
\end{equation}
and therefore
\begin{equation}
 \Pi_\Theta=I-G_\Theta(D_TG_\Theta)^{-1}D_T.
 \label{eq:intro_projector}
\end{equation}
This is the central numerical structure of CCOP.  The terminal operator fixes
the constraint that is actually accepted; the placed pressure operator fixes
the correction range.  The oblique projector in
\eqref{eq:intro_projector} is the Schur-complement realization of this
configuration assignment.

This distinction exposes a geometric degree of freedom while preserving the
pressure-space dimension.  We restrict the action to one real material configuration on the
affine labelled-position path
\[
 x_\Theta=(1-\Theta)x_n+\Theta x_T.
\]
The parameter $\Theta$ assigns an operator to a material geometry.  Temporal
state advancement and pressure-coefficient staging remain separate choices.
Midpoint action therefore produces the mixed response $D_TG_{1/2}$; a
common-stage midpoint projection would instead use $D_{1/2}G_{1/2}$.  The two
responses coincide under additional stage-consistency conditions.
Changing $\Theta$ therefore changes
$F_\Theta$, $J_\Theta$, $G_\Theta$, $D_TG_\Theta$, and the correction range,
while retaining the same scalar pressure unknowns.  The midpoint choice is
especially useful: it retains one pressure stage and, under stage-consistent
conditions, cancels the leading configuration-induced velocity defect.  The
numerical comparison therefore isolates configuration assignment within a
fixed spatial discretization.

The central contribution is this configuration assignment and its computable
consequences.  The analysis
links relative deformation to cross-response solvability, correction-range
variation to projector obliqueness, and represented pressure work to an exact
finite-step work--energy identity.  A fixed-space local-branch expansion of the actual CCOP equations
identifies the surviving projected variation of the pressure-action direction
as the leading local velocity defect.  The
discretization realizes these relations with a square
pressure-test response, coupled terminal geometry, and full-operator transport.
It distinguishes exact represented closure from consistency with the continuum
divergence and from finite-step volume preservation.

The numerical evidence follows this distinction.  Affine flows isolate target
and action effects, then measure their accumulation against a
same-discretization ISPH workflow.  A differential vortex exposes the spatial
effect of nonaffine shear: differentiation amplifies boundary-action error,
whereas manufactured pressure solves test how the inverse response transmits
that error to the corrected action.  The comparisons use shared spatial
operators within each case.  Case~2 uses its documented native-pressure
convention, while Cases~1 and~3 use the cofactor action.

The remainder is organized as follows.  \Cref{sec:configuration_assignment}
assigns the constraint and pressure action to their respective configurations
and derives the conservative midpoint choice.  \Cref{sec:ccop_analysis} then
derives the resulting cross-configuration response and its Schur projector.
\Cref{sec:meshless_discretization} describes the meshless realization and
full-operator transport.  The three numerical cases are presented in
\cref{sec:result1,sec:result2,sec:result3}, followed by the combined discussion
and conclusions.

\section{Configuration assignment of the pressure projection}
\label{sec:configuration_assignment}

The finite-step projection assigns two configuration-dependent operators.  The
accepted volume state determines where incompressibility is tested, while the
step-integrated pressure impulse determines where the multiplier acts.  We
derive these assignments in that order and then combine them into the pressure
response underlying the oblique projection.

\subsection{Terminal constraint and the fixed-configuration reference}

Let $\M^n\subset\R^d$ be the material support at the beginning of a step and
let
\begin{equation}
  \Phi_\alpha:\M^n\rightarrow\M^\alpha,
  \qquad F_\alpha=\nabla_X\Phi_\alpha,
  \qquad J_\alpha=\det F_\alpha
  \label{eq:material_maps}
\end{equation}
denote an admissible configuration in the step.  We assume
\begin{equation}
  0<J_{\min}\le J_\alpha\le J_{\max},
  \qquad
  \norm[L^\infty]{F_\alpha}+\norm[L^\infty]{F_\alpha^{-1}}
  \le C_F.
  \label{eq:admissible_geometry}
\end{equation}
Velocities are measured in the material mass inner product
\begin{equation}
  \ip[m]{u}{v}=\int_{\M^n}\rho_0 u\cdot v\,dX,
  \qquad 0<\rho_{\min}\le\rho_0\le\rho_{\max}.
  \label{eq:material_mass_inner_product}
\end{equation}
The pressure space is either $H^1(\M^n)/\R$ with a mean-zero representative,
or $H^1_{\Gamma_D^n}(\M^n)$ when a pressure trace is prescribed.

For $q$ in this gauge-fixed space, define the pulled-back weak divergence by
\begin{equation}
 \ip[Q',Q]{\D_\alpha v}{q}
 =-\int_{\M^n}J_\alpha F_\alpha^{-1}v\cdot\nabla_Xq\,dX
   +B_\alpha(v,q),
 \label{eq:weak_divergence_continuum}
\end{equation}
where
\begin{equation}
 B_\alpha(v,q)=\int_{\partial\M^n}
 q\,(J_\alpha F_\alpha^{-1}v)\cdot N\,dS.
 \label{eq:boundary_form}
\end{equation}
For smooth fields, the associated physical divergence is
\begin{equation}
  \D_\alpha^{\rm str}v
  =J_\alpha^{-1}\Div_X(J_\alpha F_\alpha^{-1}v).
  \label{eq:strong_piola_divergence}
\end{equation}
The material-cofactor pressure action is
\begin{equation}
  \G_\alpha q
  =\rho_0^{-1}J_\alpha F_\alpha^{-T}\nabla_Xq
  =\rho_0^{-1}\cof(F_\alpha)\nabla_Xq.
  \label{eq:continuum_pressure_action}
\end{equation}
If the boundary term in \eqref{eq:weak_divergence_continuum} vanishes in the
actual trial and test spaces, then
\begin{equation}
  \ip[m]{v}{\G_\alpha q}
  =-\ip[Q',Q]{\D_\alpha v}{q}.
  \label{eq:continuum_adjointness}
\end{equation}
Equation~\eqref{eq:continuum_adjointness} is the fixed-configuration Hodge
reference.  A finite Lagrangian step introduces a different question because
the accepted constraint and the pressure action need not belong to the same
configuration.

\subsubsection{Finite-volume state and its terminal tangent}
\label{sec:state_tangent_bridge}

Because $\M^n$ is the current reference, local finite-step volume preservation
is the state equation
\begin{equation}
 \mathcal C_T^J(\Phi_T)=J_T-1=0,
 \qquad
 \mathcal C_T^{\log}(\Phi_T)=\log J_T=0.
 \label{eq:finite_volume_state_residuals}
\end{equation}
The two residuals have the same zero set on the admissible family.  Their
linearization supplies the geometric meaning of terminal divergence.

\begin{proposition}[State--tangent identity]
\label{prop:state_tangent_identity}
Let $\Phi_T$ satisfy \eqref{eq:admissible_geometry}, and perturb it by a
material-labelled physical vector field $w$ as
$\Phi_{T,\varepsilon}=\Phi_T+\varepsilon w$.  Then
\begin{equation}
 D\mathcal C_T^J(\Phi_T)[w]=J_T\D_T^{\rm str}w,
 \qquad
 D\mathcal C_T^{\log}(\Phi_T)[w]=\D_T^{\rm str}w.
 \label{eq:state_tangent_variations}
\end{equation}
Consequently, if $J_T=1$, $\kernel\D_T^{\rm str}$ is the tangent space of the
volume-preserving configuration manifold at $\Phi_T$.
\end{proposition}

\begin{proof}
Jacobi's formula gives
\[
 \left.\frac{d}{d\varepsilon}\det(F_T+\varepsilon\Grad_Xw)
 \right|_{\varepsilon=0}
 =J_T\tr(F_T^{-1}\Grad_Xw).
\]
The Piola identity reduces
$J_T^{-1}\Div_X(J_TF_T^{-1}w)$ to the same trace.  Division by $J_T$ gives
the logarithmic variation.
\end{proof}

The spatial method realizes this continuum tangent in a finite pressure-test
space.  Its consistency with $D_T^{\rm str}$ is defined in
\cref{sec:discrete_operator_consistency}.  The configuration assignment itself
follows from comparing the accepted terminal state with a lagged constraint.

\subsubsection{Why the constraint is terminal}

Let $\alpha<T$ and let
$\Psi_{\alpha\to T}=\Phi_T\circ\Phi_\alpha^{-1}$ denote the relative map.
Written in $\M^\alpha$ coordinates, its terminal divergence is
\begin{equation}
 \D_T^{(\alpha)}v
 =J_{\alpha\to T}^{-1}\nabla_\alpha\cdot
 \left(J_{\alpha\to T}F_{\alpha\to T}^{-1}v\right).
 \label{eq:terminal_divergence_relative}
\end{equation}
The source discrepancy introduced by a nonterminal anchor is
\begin{equation}
  \mathcal S_{\alpha\to T}[v]
  =\dt^{-1}\left(\D_T^{(\alpha)}v-\D_\alpha v\right).
  \label{eq:source_defect_definition}
\end{equation}
Suppose
$\Psi_{\alpha\to T}(x)=x+\delta w_\alpha(x)+O(\delta^2)$ in $C^2$.
Expanding the Piola factors gives
\begin{equation}
 \D_T^{(\alpha)}v
 =\D_\alpha v+\delta\mathfrak M_\alpha(v,w_\alpha)+O(\delta^2),
 \label{eq:source_defect_expansion}
\end{equation}
where
\begin{equation}
 \mathfrak M_\alpha(v,w)
 =v\cdot\nabla_\alpha(\nabla_\alpha\cdot w)
  -\nabla_\alpha\cdot\bigl((\nabla_\alpha w)v\bigr).
 \label{eq:source_defect_leading_term}
\end{equation}
For smooth fields, commuting derivatives and expanding the second divergence
give the more transparent form
\begin{equation}
 \mathfrak M_\alpha(v,w)
 =-\tr\!\left((\nabla_\alpha w)(\nabla_\alpha v)\right).
 \label{eq:source_defect_trace_form}
\end{equation}
Thus the leading mismatch is a deformation--velocity-gradient contraction;
it need not vanish when both fields are divergence free.
Consequently, an exact solve of
\begin{equation}
  \D_\alpha\left(u^\star-\dt\G_\Th p_\alpha\right)=0
  \label{eq:nonterminal_anchor_equation}
\end{equation}
generally leaves
\begin{equation}
 \D_T^{(\alpha)}u_\alpha^T
 =\delta\mathfrak M_\alpha(u_\alpha^T,w_\alpha)+O(\delta^2).
 \label{eq:terminal_residual_after_lagged_solve}
\end{equation}
Thus accurate closure of the wrong anchor is not terminal closure.  Since the
accepted velocity is passed to the next step on $\M^T$, the constraint
assignment is
\begin{equation}
  \boxed{\alpha_c=T.}
  \label{eq:terminal_assignment}
\end{equation}

\subsection{Physical selection of the pressure-action configuration}
\label{sec:finite_step_action_work}

The pressure contribution to material momentum is an interval impulse,
\begin{equation}
  I_p=\int_{t_n}^{t_{n+1}}\G(t)p(t)\,dt.
  \label{eq:exact_pressure_impulse}
\end{equation}
A one-stage correction replaces this field-valued integral by
\begin{equation}
  I_p\approx\dt\G_\Th p_\Th.
  \label{eq:placed_pressure_impulse}
\end{equation}
We represent the action geometry on the labelled-position path
\begin{equation}
  x_\Th=x_n+\Th(x_T-x_n),\qquad 0\le\Th\le1.
  \label{eq:linear_action_path}
\end{equation}
The action configuration defines a one-stage representative of the pressure
impulse, whereas the terminal configuration defines the admissibility target.
Three choices are distinct: the temporal stage of the state integrator, the
stage represented by the pressure coefficient, and the material configuration
carrying the pressure operator.  The parameter $\Th$ controls the third choice.
It is geometric rather than merely temporal.

The action placement is selected by physical pressure-impulse and work
accuracy, rather than introduced as a tuning parameter.  To separate this
selection principle from an algebraic identity, write
\begin{equation}
  u_T=u^\star-\dt\G_\Th p_\Th,
  \qquad \D_Tu_T=0.
  \label{eq:terminal_action_step_continuum}
\end{equation}
Every admissible placement satisfies an exact work identity with the action it
represents.  This identity does not by itself determine which represented
action best approximates the physical pressure work.

\begin{proposition}[Discrete pressure work--energy identity]
\label{prop:discrete_pressure_work_energy}
Let $g_\Th=\G_\Th p_\Th$, let
$u_T=u^\star-\dt g_\Th$, and define
$E_m(u)=\half\norm[m]{u}^2$ using the same material mass inner product at both
ends of the correction.  Then, for every admissible action placement,
\begin{equation}
 \Delta E_p(\Th)
 =E_m(u_T)-E_m(u^\star)
 =-\dt\ip[m]{\frac{u^\star+u_T}{2}}{g_\Th}
 =-\dt\ip[m]{u^\star}{g_\Th}
  +\frac{\dt^2}{2}\norm[m]{g_\Th}^2.
 \label{eq:discrete_pressure_work_energy}
\end{equation}
Thus the discrete pressure work
$W_{p,\Th}^{\dt}:=-\dt\ip[m]{(u^\star+u_T)/2}{g_\Th}$ exactly accounts for
the pressure-induced kinetic-energy change.  This identity requires neither
adjointness of $D_T$ and $G_\Th$ nor symmetry of $D_TG_\Th$.
\end{proposition}

\begin{proof}
The polarization identity in the material mass inner product gives
\[
 E_m(u_T)-E_m(u^\star)
 =\ip[m]{\frac{u^\star+u_T}{2}}{u_T-u^\star}.
\]
Substitution of $u_T-u^\star=-\dt g_\Th$ proves the first equality in
\eqref{eq:discrete_pressure_work_energy}; substituting
$u_T=u^\star-\dt g_\Th$ gives the second.
\end{proof}

We call \eqref{eq:discrete_pressure_work_energy} exact discrete work
consistency with the represented pressure action.  Physical work accuracy
instead compares $W_{p,\Th}^{\dt}$ with
\begin{equation}
 W_p^{\rm ex}
 =-\int_{t_n}^{t_{n+1}}\ip[m]{u(t)}{\G(t)p(t)}\,dt.
 \label{eq:exact_continuous_pressure_work}
\end{equation}
This approximation depends on the geometric action placement, pressure-stage
consistency, the nonpressure predictor, and any force or viscous splitting.
Accordingly, $D_T$ selects the admissible terminal state, whereas
$G_\Th p_\Th$ selects the pressure direction and the work represented by the
correction.  Terminal closure alone therefore fixes neither the pressure
impulse nor its kinetic work.

\subsection{Why midpoint action is selected in the present conservative setting}

The temporal consequence of action placement follows from a one-point
quadrature expansion once the coefficient stage and geometric path satisfy
their own consistency conditions.

\begin{proposition}[Defect of the placed pressure impulse]
\label{prop:placed_impulse_order}
Let $g(t)=\G(t)p(t)$ be twice continuously differentiable on one smooth
semidiscrete step, and define the action-stage error in the momentum norm by
\begin{equation}
  \eta_\Th^h=g_\Th^h-g(t_n+\Th\dt).
\label{eq:second_order_stage_consistency}
\end{equation}
Then
\begin{equation}
 \int_{t_n}^{t_{n+1}}g(t)\,dt-\dt g_\Th^h
 =\left(\frac12-\Th\right)\dt^2\dot g(t_n)
  -\dt\eta_\Th^h+O(\dt^3).
 \label{eq:placed_impulse_defect}
\end{equation}
Consequently, if $\eta_{1/2}^h=O(\dt^2)$, midpoint placement has a
third-order local pressure-impulse defect.  For any fixed
$\Th\ne1/2$, the quadrature term is generically $O(\dt^2)$; an
$O(\dt)$ endpoint stage error has the same local order.  These are local
statements.  Global temporal rates of the complete nonlinear scheme are
assessed separately in the numerical experiments.
\end{proposition}

\begin{proof}
Taylor expansion about $t_n$ gives
\[
 \int_{t_n}^{t_{n+1}}g(t)\,dt
 =\dt g(t_n)+\frac{\dt^2}{2}\dot g(t_n)+O(\dt^3)
\]
and
\[
 \dt g(t_n+\Th\dt)
 =\dt g(t_n)+\Th\dt^2\dot g(t_n)+O(\dt^3).
\]
Subtracting $\dt\eta_\Th^h$ proves
\eqref{eq:placed_impulse_defect}.  The stated local orders follow by
cancellation at $\Th=1/2$.
\end{proof}

The proposition separates impulse accuracy into quadrature, pressure-stage,
and geometric-placement contributions.  It therefore determines the accuracy
of the action entering the exact work identity, rather than the validity of
that identity.

\subsubsection{Coupled local branch and configuration-placement defect}

The actual CCOP residual is degenerate at zero step size: if $D_nu_n=0$, then
$\mathcal R(p,0)=0$ for every pressure coefficient.  The following rescaling
first establishes the numerical pressure branch needed for a local expansion.
Indeed, because $D_px_T=O_h(\tau^2)$,
\begin{equation}
 D_p\mathcal R(p,\tau)=-\tau D_nG_n+O_h(\tau^2),
 \label{eq:degenerate_pressure_jacobian}
\end{equation}
so the unscaled residual does not satisfy the hypotheses of the ordinary
implicit-function theorem at $\tau=0$.

\begin{lemma}[Rescaled local CCOP branch]
\label{lem:regularized_local_ccop_branch}
Fix a finite-dimensional velocity space $V_h$, a pressure space $Q_{h,0}$
after imposing the applicable gauge or homogeneous boundary condition, and an
equal-dimensional constraint space $Z_h$ (realized below as the corresponding
pressure-test dual), together with fixed pressure labels.  Consider
the smooth semidiscrete constrained system
\begin{equation}
 \dot x=u,\qquad \dot u=f(x)-G(x)p,\qquad D(x)u=0.
 \label{eq:semidiscrete_constrained_flow}
\end{equation}
Assume $f\in C^2$, $D\in C^3$, $G\in C^2$, $D_nu_n=0$, and
$L_n=D_nG_n:Q_{h,0}\to Z_h$ is bijective.  Let
$u^\star(\tau)$ be $C^3$, let $a^\star(\tau)$ be $C^2$, and write
\begin{equation}
 u^\star(0)=u_n,
 \qquad
 x^\star(\tau)=x_n+\tau u_n+\tau^2a^\star(\tau).
 \label{eq:general_predictor_scaling}
\end{equation}
For fixed $\Th\in[0,1]$, set
\begin{equation}
 x_T=x_n+\tau u_n+\tau^2\xi,
 \qquad
 x_\Th=x_n+\Th(x_T-x_n).
 \label{eq:regularized_terminal_scaling}
\end{equation}
Then the CCOP equations have a unique local $C^2$ branch
$(\xi_\Th(\tau),p_\Th(\tau))$ near $\tau=0$.  Its value at zero is
\begin{align}
 p_0&=L_n^{-1}\left(D_n'[u_n]u_n+D_n\dot u^\star(0)\right),
 \label{eq:regularized_pressure_limit}\\
 \xi_0&=a^\star(0)-\gamma G_np_0.
 \label{eq:regularized_position_limit}
\end{align}
If $\dot u^\star(0)=f_n$, then $p_0=p_n$.  If, in addition,
$a^\star(0)=f_n/2$ and $\gamma=1/2$, then
$\xi_0=\dot u_n/2$.
\end{lemma}

\begin{proof}
For $\tau\ne0$, the position and terminal-constraint equations are equivalent
to
\begin{align}
 \Phi_1(\xi,p;\tau)
 &=\xi-a^\star(\tau)+\gamma G(x_\Th)p=0,
 \label{eq:regularized_ccop_position}\\
 \Phi_2(\xi,p;\tau)
 &=\tau^{-1}D(x_T)u^\star(\tau)-D(x_T)G(x_\Th)p=0.
 \label{eq:regularized_ccop_constraint}
\end{align}
For fixed $\xi$, set
$r(\tau,\xi)=D(x_n+\tau u_n+\tau^2\xi)u^\star(\tau)$.
Since $r(0,\xi)=D_nu_n=0$, Hadamard's lemma gives the $C^2$ extension
\begin{equation}
 \frac{r(\tau,\xi)}{\tau}
 =\int_0^1\partial_\tau r(s\tau,\xi)\,ds,
 \qquad
 \left.\frac{r(\tau,\xi)}{\tau}\right|_{\tau=0}
 =D_n'[u_n]u_n+D_n\dot u^\star(0).
 \label{eq:hadamard_constraint_extension}
\end{equation}
Hence \eqref{eq:regularized_ccop_position}--
\eqref{eq:regularized_ccop_constraint} extend to a $C^2$ map at zero.  At the
root defined by \eqref{eq:regularized_pressure_limit}--
\eqref{eq:regularized_position_limit}, its Jacobian is
\begin{equation}
 D_{(\xi,p)}\Phi(\xi_0,p_0;0)
 =\begin{bmatrix}
 I & \gamma G_n\\
 0 & -D_nG_n
 \end{bmatrix}.
 \label{eq:regularized_ccop_jacobian}
\end{equation}
This matrix is invertible precisely when $L_n$ is bijective on the stated
spaces, so the implicit-function theorem supplies the local branch.  Finally,
differentiating $D(x)u=0$ along
\eqref{eq:semidiscrete_constrained_flow} gives
\begin{equation*}
 D_nG_np_n=D_n'[u_n]u_n+D_nf_n.
\end{equation*}
The last two conclusions follow from uniqueness and
\eqref{eq:regularized_position_limit}.
\end{proof}

Here and below, primes denote Fr\'echet derivatives with respect to the
configuration.  The limiting pressure is independent of $\Th$ because every action
configuration reduces to $x_n$ at $\tau=0$; placement enters the next
coefficient of the branch.  For the implemented fixed-stencil, fixed-label
operators, the determinant, inverse, and cofactor formulas are analytic on the
open set $j_Tj_\Th\ne0$, which is monitored by the raw-determinant acceptance
test.  The expansion starts from exact represented closure.  An incoming
residual $D_nu_n=\epsilon$ enters the rescaled equation as
$\epsilon/\tau$ and is tracked as a numerical perturbation.  The lemma
characterizes the local solution branch; convergence of the block Picard map
is established separately in \cref{thm:local_coupled_closure}.

\begin{proposition}[Local configuration-placement defect]
\label{prop:actual_ccop_local_defect}
Assume the hypotheses of
\cref{lem:regularized_local_ccop_branch}, let the exact solution and pressure
of \eqref{eq:semidiscrete_constrained_flow} be sufficiently smooth, and take
$\gamma=1/2$.  Suppose the configuration-force predictor satisfies
\begin{align}
 u^\star(\tau)
 &=u_n+\tau f_n+\frac{\tau^2}{2}f_n'[u_n]+O_h(\tau^3),
 \label{eq:second_order_force_velocity_predictor}\\
 x^\star(\tau)
 &=x_n+\tau u_n+\frac{\tau^2}{2}f_n+O_h(\tau^3).
 \label{eq:second_order_force_position_predictor}
\end{align}
Then, in any fixed equivalent finite-dimensional norm,
\begin{equation}
 u_{T,\Th}^h-u_h(t_n+\tau)
 =\left(\frac12-\Th\right)\tau^2
   \Pi_n G_n'[u_n]p_n+O_h(\tau^3),
 \qquad
 \Pi_n=I-G_n(D_nG_n)^{-1}D_n.
 \label{eq:actual_ccop_local_velocity_defect}
\end{equation}
If the stated smoothness and inverse bounds are uniform over a spatial family,
the remainder constant is uniform in $h$.  Otherwise
\eqref{eq:actual_ccop_local_velocity_defect} is a fixed-space local result and
must be combined with a separate spatial-consistency estimate.
\end{proposition}

\begin{proof}
The lemma gives
$p_\Th^h=p_n+\tau\pi_1+O_h(\tau^2)$.  The predictor and trapezoidal
kinematics give
\begin{equation*}
 x_T=x_n+\tau u_n+\frac{\tau^2}{2}(f_n-G_np_n)+O_h(\tau^3),
\end{equation*}
so the numerical and exact terminal positions agree through $O_h(\tau^2)$.
Moreover,
\begin{equation*}
 G(x_\Th)p_\Th^h
 =G_np_n+\tau\bigl(G_n\pi_1+\Th G_n'[u_n]p_n\bigr)
  +O_h(\tau^2).
\end{equation*}
The exact semidiscrete velocity satisfies
\begin{equation*}
 u_h(t_n+\tau)=u_n+\tau(f_n-G_np_n)
 +\frac{\tau^2}{2}
 \bigl(f_n'[u_n]-G_n'[u_n]p_n-G_n\dot p_n\bigr)
 +O_h(\tau^3).
\end{equation*}
Subtracting the exact and numerical terminal constraints shows that the
coefficient of $\tau^2$ in the velocity difference lies in $\kernel D_n$.
With $\delta p=\pi_1-\dot p_n/2$, this condition is
\begin{equation*}
 D_nG_n\delta p
 =-\left(\Th-\frac12\right)D_nG_n'[u_n]p_n.
\end{equation*}
Solving on $Q_{h,0}$ and substituting into the velocity expansion gives
\eqref{eq:actual_ccop_local_velocity_defect}.
\end{proof}

The proof also identifies the multiplier stage:
\begin{equation}
 p_\Th^h
 =p_n+\frac{\tau}{2}\dot p_n
  -\left(\Th-\frac12\right)\tau
   (D_nG_n)^{-1}D_nG_n'[u_n]p_n
  +O_h(\tau^2).
 \label{eq:actual_ccop_multiplier_stage}
\end{equation}
Thus the multiplier equation centers the pressure coefficient through first
order: the midpoint coefficient represents both the instantaneous midpoint
pressure and the step-average pressure to the stated accuracy, whereas generic
endpoint coefficients differ from their endpoint samples at $O(\tau)$.  This
coefficient centering does not center the geometric pressure operator.  The
latter is the independent configuration choice controlled by $\Th$.

For $f(x)=-A_fx$, expansion of the implicit-midpoint predictor
\eqref{eq:conservative_force_predictor} gives
\begin{align*}
 u^\star
 &=u_n-\tau A_fx_n-\frac{\tau^2}{2}A_fu_n+O_h(\tau^3)
  =u_n+\tau f_n+\frac{\tau^2}{2}f_n'[u_n]+O_h(\tau^3),\\
 x^\star
 &=x_n+\tau u_n-\frac{\tau^2}{2}A_fx_n+O_h(\tau^3)
  =x_n+\tau u_n+\frac{\tau^2}{2}f_n+O_h(\tau^3).
\end{align*}
Thus \eqref{eq:second_order_force_velocity_predictor}--
\eqref{eq:second_order_force_position_predictor} hold, and
\cref{prop:actual_ccop_local_defect} covers force-free Case~1 and the
conservative-force architecture used in Cases~2 and~3.

The leading term has a direct geometric interpretation: only the component of
the changing pressure-action direction that survives the current constrained
projector contributes to the local velocity defect.  Thus midpoint placement
cancels the $O(\dt^2)$ configuration term for the actual CCOP solve under
\eqref{eq:second_order_force_velocity_predictor}--
\eqref{eq:second_order_force_position_predictor}.  Relative to
a continuum solution, the corresponding stage estimate must retain the
spatial floor, for example
\begin{equation}
 \norm{\eta_{1/2,h}}
 \le C_t\dt^2+C_sh^r,
 \label{eq:space_time_stage_floor}
\end{equation}
where $r$ is established by the spatial action audit.  Velocity-dependent
viscous terms, nonmatching force splittings, and nonhomogeneous-boundary stages
require separate consistency arguments.

Proposition~\ref{prop:placed_impulse_order} therefore gives a pressure-action
design rule under stage consistency.  A second-order midpoint action gives a
third-order local impulse defect, whereas a first-order endpoint stage gives a
second-order local defect.  Proposition~\ref{prop:actual_ccop_local_defect}
shows the corresponding cancellation in the fully coupled velocity update.
Varying $\Th$ changes the action geometry within the same scalar pressure
stage.

For compatible same-configuration endpoint pairs, the work identity also gives
a sign structure that identifies an interior work-neutral action.

\begin{proposition}[Endpoint work bias]
\label{prop:endpoint_work_bias}
Assume compatible boundary terms at $n$ and $T$, exact endpoint solves, and
$\D_nu^\star=0$.  Then
\begin{equation}
  \Delta E_p(n)=\frac{\dt^2}{2}\norm[m]{\G_np_n}^2\ge0,
  \qquad
  \Delta E_p(T)=-\frac{\dt^2}{2}\norm[m]{\G_Tp_T}^2\le0.
  \label{eq:endpoint_work_signs}
\end{equation}
\end{proposition}

\begin{proof}
At the old endpoint,
$\ip[m]{u^\star}{\G_np_n}=-\ip{\D_nu^\star}{p_n}=0$.
At the terminal endpoint, the pressure equation and
\eqref{eq:continuum_adjointness} give
$\ip[m]{u^\star}{\G_Tp_T}=\dt\norm[m]{\G_Tp_T}^2$.
Substitution into \eqref{eq:discrete_pressure_work_energy} proves the
claim.
\end{proof}

Whenever the endpoint values bracket zero, continuity supplies a work-neutral
interior placement.  The impulse expansion and the coupled local defect select
the midpoint to leading order for the stage-consistent conservative predictor
used here.  Case~1 independently measures the work-neutral root and its
convergence to the midpoint.  The general configuration assignment and the
choice used for the principal conservative comparisons are therefore
\begin{equation}
  \boxed{(\alpha_c,\alpha_a)=(T,\Th),
  \qquad \Th=\frac12\ \text{for the present conservative setting}.}
  \label{eq:configuration_assignment_summary}
\end{equation}

These two assignments determine the pressure response.  Terminal admissibility
is tested by $D_T$, while the pressure action is represented by $G_\Th$.
Consequently, the frozen pressure derivative of the terminal residual is
\begin{equation}
 D_p\!\left[\D_T
 (u^\star-\dt\G_\Th p)\right][\delta p]
 =-\dt\D_T\G_\Th\delta p.
 \label{eq:cross_response_as_tangent_pressure_block}
\end{equation}
Thus $D_TG_\Th$ is the pressure-to-terminal-constraint response selected by the
two physical roles.  Its inverse generates the configuration-dependent oblique
projection developed next.  In the coupled step, the pressure dependence of
$x_T$ adds geometry derivatives around this frozen principal response.

\section{Configuration-consistent oblique projection}
\label{sec:ccop_analysis}

\subsection{Action path and nonlinear terminal state}

At the beginning of a step, $x_n$ is the current labelled configuration.  The
nonpressure predictor supplies $u^\star$ and $x^\star$.  The action path
\eqref{eq:linear_action_path} parameterizes admissible pressure-operator
geometries independently of the state time integrator.  Hence $\Th=1/2$ defines the
mixed-configuration response $D_TG_{1/2}$, while a common-stage midpoint
method would use $D_{1/2}G_{1/2}$.  Every placement retains a single scalar
pressure/PPE solve.  With the current spatial gradient $\Grad_n$,
\begin{align}
 H_T&=I+\Grad_n(x_T-x_n), & j_T&=\det H_T,\label{eq:relative_terminal_F}\\
 H_\Th&=I+\Th(H_T-I), & j_\Th&=\det H_\Th.
 \label{eq:relative_action_F}
\end{align}
The action state is the numerical path state selected by
\eqref{eq:linear_action_path}.  Its raw determinant is monitored for
admissibility independently from the volume-preserving terminal state.

For the pressure part of the position update we use $\gamma=1/2$ and solve
\begin{align}
 u_T&=u^\star-\dt G_\Th[x_T]p,\label{eq:nonlinear_velocity}\\
 x_T&=x^\star-\gamma\dt^2G_\Th[x_T]p,\label{eq:nonlinear_position}\\
 D_T[x_T]u_T&=0.\label{eq:nonlinear_terminal_constraint}
\end{align}
Eliminating $u_T$ and the implicit kinematic relation defines a pressure-only
residual,
\begin{equation}
 \mathcal R(p)=D_T[x_T(p)]
 \left(u^\star-\dt G_\Th[x_T(p)]p\right),
 \quad
 x_T(p)=x^\star-\gamma\dt^2G_\Th[x_T(p)]p.
 \label{eq:reduced_nonlinear_residual}
\end{equation}
Trapezoidal particle kinematics fixes $\gamma=1/2$, independently of the
work-neutral placement.  In the force-free case
$u^\star=u_n$ and $x^\star=x_n+\dt u_n$, so
\begin{equation}
 x_T=x_n+\frac{\dt}{2}(u_n+u_T)
     =x^\star-\frac{\dt^2}{2}G_\Th p.
 \label{eq:gamma_trapezoidal_origin}
\end{equation}
The conservative-force predictor used later builds the same trapezoidal
kinematic relation into $x^\star$ and satisfies
\eqref{eq:second_order_force_velocity_predictor}--
\eqref{eq:second_order_force_position_predictor}.  Velocity-dependent viscous
or split operators are outside the local result.  The projector introduced
below is the frozen linear algebra embedded in this nonlinear geometry
coupling.

\subsection{Weak placed problem and relative-deformation solvability}

Freeze admissible maps $F_T$ and $F_\Th$.  Using
\eqref{eq:weak_divergence_continuum}--\eqref{eq:continuum_pressure_action} and
compatible boundary terms, the pressure equation is the variational problem
\begin{equation}
  a_{T,\Th}(p,q)=\dt^{-1}b_T(u^\star,q)
  \qquad\text{for all }q\in Q,
  \label{eq:weak_placed_problem}
\end{equation}
with
\begin{align}
 a_{T,\Th}(p,q)
 &=\int_{\M^n}\nabla_Xq\cdot A_{T,\Th}\nabla_Xp\,dX,
 \label{eq:placed_bilinear}\\
 A_{T,\Th}
 &=\rho_0^{-1}J_TJ_\Th F_T^{-1}F_\Th^{-T},
 \label{eq:relative_tensor}\\
 b_T(v,q)
 &=\int_{\M^n}J_TF_T^{-1}v\cdot\nabla_Xq\,dX.
 \label{eq:terminal_load}
\end{align}

\begin{theorem}[Frozen placed solvability]
\label{thm:placed_solvability}
Suppose \eqref{eq:admissible_geometry} holds and there is $c_a>0$ such that
\begin{equation}
  \xi\cdot\sym(A_{T,\Th}(X))\xi\ge c_a|\xi|^2
  \label{eq:coercivity_assumption}
\end{equation}
for almost every $X$ and every $\xi\in\R^d$.  Then
\eqref{eq:weak_placed_problem} has a unique solution in the gauge-fixed
pressure space.
\end{theorem}

\begin{proof}
The geometry and density bounds make $a_{T,\Th}$ and $b_T$ bounded.  The
lower bound \eqref{eq:coercivity_assumption}, together with the quotient or
Dirichlet Poincar\'e inequality, makes $a_{T,\Th}$ coercive.  The result follows
from the Lax--Milgram theorem.
\end{proof}

The condition has a direct relative-deformation interpretation.  If
$z=F_T^{-T}\xi$, then
\begin{equation}
 \xi^TF_T^{-1}F_\Th^{-T}\xi
 =z^T\sym(F_TF_\Th^{-1})z.
 \label{eq:relative_deformation_rewrite}
\end{equation}
Hence positivity of
$\sym(F_TF_\Th^{-1})$ is sufficient after the geometry bounds are included.
In particular,
\begin{equation}
 \norm[2]{F_TF_\Th^{-1}-I}\le\delta<1
 \quad\Longrightarrow\quad
 \lambda_{\min}\!\left(\sym(F_TF_\Th^{-1})\right)\ge1-\delta.
 \label{eq:small_relative_deformation_condition}
\end{equation}
This statement acts directly on the actual relative tensor and requires no
commutation with a perturbation factor.

For the linear action path this condition can be checked before a solve.  Let
$E=H_T-I$ and suppose $\norm[2]{E}=\varepsilon<1$.  Since
\begin{equation}
 H_TH_\Th^{-1}-I
 =(1-\Th)E(I+\Th E)^{-1},
 \label{eq:relative_path_exact_identity}
\end{equation}
the Neumann bound gives
\begin{equation}
 \norm[2]{H_TH_\Th^{-1}-I}
 \le \frac{(1-\Th)\varepsilon}{1-\Th\varepsilon}<1.
 \label{eq:relative_path_explicit_bound}
\end{equation}
If $\varepsilon\le \dt\norm[L^\infty]{\nabla u}+C\dt^2$, this is an explicit
small-deformation step condition.  It is sufficient, not necessary: the
experiments also measure the eigenvalue margin directly.

\subsection{Frozen projection algebra}

Let $L_{T,\Th}=D_TG_\Th$ be invertible on the selected pressure space.  Define
\begin{equation}
  \Pi_\Th=I-G_\Th L_{T,\Th}^{-1}D_T.
  \label{eq:frozen_ccop_operator}
\end{equation}

\begin{proposition}[Exact frozen projection]
\label{prop:frozen_projection}
The operator \eqref{eq:frozen_ccop_operator} satisfies
\begin{equation}
 D_T\Pi_\Th=0,
 \qquad
 \Pi_\Th^2=\Pi_\Th,
 \qquad
 \Pi_\Th G_\Th=0,
 \label{eq:frozen_projection_identities}
\end{equation}
and
\begin{equation}
 \range(\Pi_\Th)=\kernel(D_T),
 \qquad
 \kernel(\Pi_\Th)=\range(G_\Th).
 \label{eq:frozen_projection_spaces}
\end{equation}
\end{proposition}

\begin{proof}
Writing $Q_\Th=G_\Th L_{T,\Th}^{-1}D_T$ gives
$Q_\Th^2=Q_\Th$, $D_TQ_\Th=D_T$, and
$Q_\Th G_\Th=G_\Th$.  The identities and the range/kernel statements follow
from $\Pi_\Th=I-Q_\Th$.
\end{proof}

At the continuum endpoint, \eqref{eq:continuum_adjointness} makes
$\range(\G_T)$ the mass-normal complement of $\kernel(\D_T)$; the ideal
$\Pi_T$ is then the terminal Hodge projector.  For $\Th\ne T$, equality of the
two action ranges would require, for every representable $p$, a scalar
$\widetilde p$ satisfying
\begin{equation}
 \nabla_X\widetilde p
 =\frac{J_\Th}{J_T}F_T^TF_\Th^{-T}\nabla_Xp.
 \label{eq:off_terminal_integrability}
\end{equation}
The right side is not generally curl free, so the off-terminal action range is
generically not the terminal normal space.

The realized discrete endpoint has a measured nonorthogonal baseline.
Accordingly, the Hodge terminology is reserved for the compatible continuum
statement.  For the realized operators we report the endpoint norm as a
baseline and use
$\norm[M_U]{\Pi_\Th-\Pi_T}$ and the principal angle between
$\range(G_\Th)$ and $\range(G_T)$ to measure the additional
configuration-induced obliqueness.

The following auxiliary construction separates that increment from any
nonorthogonality already present in the meshless pair.  It is a diagnostic
control and is not used in the pressure solve.  Define the exact terminal
mass-adjoint functional and its placed projector by
\begin{equation}
 B_T^{\rm ad}=-G_T^TM_U,
 \qquad
 \Pi_\Th^{\rm ad}
 =I-G_\Th(B_T^{\rm ad}G_\Th)^{-1}B_T^{\rm ad}.
 \label{eq:compatible_endpoint_control}
\end{equation}
The code stores the point-value representative
$D_T^{\rm ad}=M_Z^{-1}B_T^{\rm ad}$; the left Riesz factor cancels from the
projector, so this is algebraically identical to
\eqref{eq:compatible_endpoint_control}.
If $G_T$ has full column rank, then $\Pi_T^{\rm ad}$ is the $M_U$-orthogonal
projector onto $\ker(B_T^{\rm ad})$.  For $\Th\ne T$,
$\Pi_\Th^{\rm ad}$ has the same target kernel but null space
$\range(G_\Th)$ and is generally oblique.  Hence this control has an exactly
orthogonal endpoint baseline and attributes every off-endpoint departure to
the change of action range.
Indeed, at $\Th=T$ the minus signs cancel and
$\Pi_T^{\rm ad}=I-G_T(G_T^TM_UG_T)^{-1}G_T^TM_U$, the standard
$M_U$-orthogonal projector.  For general $\Th$, the algebra in
\cref{prop:frozen_projection} applies with $D_T=B_T^{\rm ad}$.

For the realized point-value pair, we quantify same-configuration discrete
Green compatibility by
\begin{equation}
 e_{{\rm Gr},{\rm full}}
 =\frac{\norm[F]{D_T^q-D_T^{\rm ad}}}
        {\norm[F]{D_T^q}},
 \qquad
 D_T^{\rm ad}=-M_Z^{-1}G_T^TM_U,
 \label{eq:full_space_green_compatibility}
\end{equation}
where $\norm[F]{\cdot}$ is the Frobenius norm.  Thus
$e_{{\rm Gr},{\rm full}}=0$ is equivalent to
$M_ZD_T^q=-G_T^TM_U$.  This full-matrix compatibility diagnostic is not a
truncation error for either $D_T^q$ or $G_T$, and it does not measure pressure
residual or projector closure.  The row-restricted strong GMLS construction
does not impose a discrete integration-by-parts identity, so a nonzero value is
compatible with accurate action and divergence on resolved fields.

Compatibility on a resolved physical subspace can differ from the full-space
quantity.  For velocity and pressure subspaces
$\mathcal V_r\subset V_h$ and $\mathcal Q_r\subset Q_h$, we therefore also use
\begin{equation}
 e_{{\rm Gr},r}
 =\sup_{\substack{0\ne v\in\mathcal V_r\\0\ne q\in\mathcal Q_r}}
 \frac{\left|v^TM_UG_Tq+q^TM_ZD_T^qv\right|}
      {\norm[M_U]{v}\,\norm[M_Q]{q}}.
 \label{eq:restricted_green_compatibility}
\end{equation}
The subspaces used in the affine audit are stated with its results.

\begin{proposition}[One-step projector proximity]
\label{prop:projector_proximity}
Suppose $G_\Th=G_T+\dt E_\Th$, the terminal response is invertible, and the
placed responses remain uniformly invertible.  Then
\begin{equation}
 \Pi_T-\Pi_\Th
 =\dt\,\Pi_T E_\Th L_{T,\Th}^{-1}D_T.
 \label{eq:projector_proximity_identity}
\end{equation}
Consequently $\norm{\Pi_T-\Pi_\Th}=O(\dt)$, and its action on a predictor
with $\norm{D_Tu^\star}=O(\dt)$ is $O(\dt^2)$.
\end{proposition}

\begin{proof}
Apply the resolvent identity to
$L_{T,\Th}=D_TG_\Th=L_{T,T}+\dt D_TE_\Th$ and collect terms:
\[
 G_\Th L_{T,\Th}^{-1}D_T-G_TL_{T,T}^{-1}D_T
 =\dt(I-G_TL_{T,T}^{-1}D_T)E_\Th L_{T,\Th}^{-1}D_T.
\]
This is \eqref{eq:projector_proximity_identity}; the norm statements follow
from the uniform inverse bounds.
\end{proof}

This proximity result concerns consistent small-step families.  Together with
\cref{prop:placed_impulse_order}, it separates two related effects: the
correction range changes by $O(\dt)$, while acting on an $O(\dt)$ pressure
source changes the one-step velocity by $O(\dt^2)$.  Midpoint placement
cancels the leading configuration term under the predictor hypotheses of
\cref{prop:actual_ccop_local_defect}, whereas terminal placement retains it.

The same range variation that produces projector obliqueness also controls the
frozen placement sensitivity of the pressure work.  The following identity
makes that connection precise without assuming adjointness.

\begin{proposition}[Frozen placement--work sensitivity]
\label{prop:frozen_placement_work_sensitivity}
Fix $D_T$ and $u^\star$, let $G_\Th$ be differentiable in $\Th$, and suppose
$L_{T,\Th}=D_TG_\Th$ is invertible throughout an interval of placements.  Set
\begin{equation}
 q_\Th=L_{T,\Th}^{-1}D_Tu^\star,
 \qquad
 u_\Th=\Pi_\Th u^\star.
 \label{eq:frozen_placement_pressure_and_velocity}
\end{equation}
Then
\begin{align}
 \partial_\Th u_\Th
 &=-\Pi_\Th(\partial_\Th G_\Th)q_\Th,
 \label{eq:frozen_velocity_placement_sensitivity}\\
 \frac{d}{d\Th}\!\left[
  \half\norm[m]{u_\Th}^2-
  \half\norm[m]{u^\star}^2\right]
 &=-\ip[m]{u_\Th}
 {\Pi_\Th(\partial_\Th G_\Th)q_\Th}.
 \label{eq:frozen_work_placement_sensitivity}
\end{align}
\end{proposition}

\begin{proof}
Differentiate $L_{T,\Th}q_\Th=D_Tu^\star$ to obtain
$\partial_\Th q_\Th=-L_{T,\Th}^{-1}D_T
(\partial_\Th G_\Th)q_\Th$.  Since
$u_\Th=u^\star-G_\Th q_\Th$, substitution gives
\eqref{eq:frozen_velocity_placement_sensitivity}; differentiation of the
kinetic energy gives \eqref{eq:frozen_work_placement_sensitivity}.
\end{proof}

Thus a common target kernel permits placement-dependent work: the component of
the changing action direction that survives $\Pi_\Th$ enters the first
variation.  Projector distance, principal angle, and work are complementary
diagnostics.  Their common finite-step source is the configuration-driven
variation of $\range(G_\Th)$.  In the fully coupled step, total differentiation
adds the dependence of $D_T$, $G_\Th$, and $x_T$ on the nonlinear accepted
geometry.

\subsection{Local nonlinear solvability}
\label{sec:conditional_coupled_solvability}

The frozen inverse is also the core of the nonlinear solve.  Uniform
invertibility on a closed admissible configuration neighborhood $\mathcal U$
leads to the following local closure criterion.  Define
\begin{equation}
 L_\Th[x]=D_T[x]G_\Th[x],
 \qquad
 b_\Th[x]=\dt^{-1}D_T[x]u^\star,
 \qquad
 \mathcal P_\Th[x]=L_\Th[x]^{-1}b_\Th[x],
 \label{eq:nonlinear_frozen_solution_map}
\end{equation}
on the fixed gauge-reduced pressure and constraint spaces.  The reduced
geometry map associated with \eqref{eq:nonlinear_position} is
\begin{equation}
 \mathcal H_\Th[x]
 =x^\star-\gamma\dt^2G_\Th[x]\mathcal P_\Th[x].
 \label{eq:nonlinear_reduced_geometry_map}
\end{equation}

\begin{proposition}[Local coupled closure criterion]
\label{thm:local_coupled_closure}
Assume on $\mathcal U$ that
$\norm{L_\Th[x]^{-1}}\le C_L$; that $L_\Th$, $b_\Th$, and $G_\Th$ are
Lipschitz with constants $L_L$, $L_b$, and $L_G$; and that
$\norm{G_\Th[x]}\le C_G$ and
$\norm{\mathcal P_\Th[x]}\le C_p$.  Set
\begin{equation}
 L_P=C_L(L_b+L_LC_p),
 \qquad
 \kappa_{\rm cpl}
 =\gamma\dt^2(L_GC_p+C_GL_P).
 \label{eq:coupled_contraction_constant}
\end{equation}
If $\mathcal H_\Th(\mathcal U)\subset\mathcal U$ and
$\kappa_{\rm cpl}<1$, then the coupled system
\eqref{eq:nonlinear_velocity}--\eqref{eq:nonlinear_terminal_constraint} has a
unique solution in $\mathcal U$.  Moreover, the frozen-pressure/block-Picard
iteration $p^{(k+1)}=\mathcal P_\Th[x^{(k)}]$,
$x^{(k+1)}=\mathcal H_\Th[x^{(k)}]$ converges to that solution.
\end{proposition}

\begin{proof}
Subtracting
$L_\Th[x]\mathcal P_\Th[x]=b_\Th[x]$ and
$L_\Th[y]\mathcal P_\Th[y]=b_\Th[y]$ gives
$\norm{\mathcal P_\Th[x]-\mathcal P_\Th[y]}
\le L_P\norm{x-y}$.  Adding and subtracting
$G_\Th[x]\mathcal P_\Th[y]$ in
\eqref{eq:nonlinear_reduced_geometry_map} then gives
$\norm{\mathcal H_\Th[x]-\mathcal H_\Th[y]}
\le\kappa_{\rm cpl}\norm{x-y}$.  Banach's fixed-point theorem supplies the
unique geometry; \eqref{eq:nonlinear_frozen_solution_map} supplies the unique
pressure, and the velocity follows from \eqref{eq:nonlinear_velocity}.
\end{proof}

For a fixed spatial discretization, the generic small-step bound scales as
$O(\dt)$ because $b_\Th$ contains $\dt^{-1}$.  No $h$-uniform contraction rate
is claimed.  In the discrete calculations, the measured Riesz gap supplies
the computable inverse factor
$C_L\le\beta_{T,\Th,h}^{-1}$ in the corresponding Riesz norms.  The iteration
histories measure convergence, and the raw Jacobians check geometric
admissibility; these measurements do not estimate all Lipschitz constants in
\eqref{eq:coupled_contraction_constant}.

\section{Meshless realization and full-operator transport}
\label{sec:meshless_discretization}

The continuous formulation determines the configuration assigned to each
operator.  The spatial realization completes the construction by forming a
stable square map between the selected pressure and constraint spaces.  This
section gives the GMLS realization used in Cases~1 and~3.  Case~2 specifies
its SPH2/MLS2 variant separately.  In both settings, endpoint adjointness is
distinct from cross-response invertibility.

\subsection{Point clouds and GMLS derivatives}

Let $\{X_i\}_{i=1}^N$ be a quasi-uniform cloud with positive material volumes
$m_i$ and boundary labels $\mathcal I_\Gamma$.  On the initial physical domain
of each run we construct degree-two GMLS first derivatives.  For a center
$X_i$, scaled monomials $P_i$ of total degree two and a compact Wendland
$C^2$ weight matrix $W_i$ give the derivative row
\begin{equation}
  w_i^{(\mu)}=W_iP_i(P_i^TW_iP_i)^{-1}\lambda_i^{(\mu)},
  \qquad \mu\in\{x,y\},
  \label{eq:gmls_row}
\end{equation}
where $\lambda_i^{(\mu)}$ evaluates the corresponding derivative of the
scaled polynomial basis at the stencil center.  The stencils used in the
pressure step have
18 neighbors.  Shifted and scaled coordinates are essential: rank is checked
against a relative singular threshold $10^{-12}$, and both the polynomial
moment defect and local condition number are recorded.  Degree-three,
32-neighbor derivatives define an independent all-particle divergence
diagnostic; they do not enter the pressure equation.

The scalar derivatives used in the pressure step are denoted
$H_x,H_y\in\R^{N\times N}$.
For interleaved velocity storage
$v=(v_{1x},v_{1y},\ldots,v_{Nx},v_{Ny})^T$, the initial all-particle
divergence is
\begin{equation}
  D_0^{\rm str}v=H_xv_x+H_yv_y.
  \label{eq:initial_strong_matrix}
\end{equation}
The same derivatives reconstruct the current vector gradient needed for
relative geometry.  Their shared polynomial moments make the affine mechanism
test exact, while the nonaffine tests directly measure spatial truncation.

\subsection{Square pressure-test spaces}

On a free surface, pressure values on $\mathcal I_\Gamma$ are prescribed and
the interior labels
\begin{equation}
  \mathcal I_q=\{1,\ldots,N\}\setminus\mathcal I_\Gamma,
  \qquad N_q=N-|\mathcal I_\Gamma|,
  \label{eq:interior_pressure_labels}
\end{equation}
are the pressure coefficients.  The initial homogeneous action and constraint
divergence are
\begin{equation}
  G_0\in\R^{2N\times N_q},
  \qquad
  D_0^q=D_0^{\rm str}[\mathcal I_q,:]
  \in\R^{N_q\times2N}.
  \label{eq:initial_square_pair}
\end{equation}
Thus $D_T^qG_\Th$ is square.  The pressure step uses this pressure-test
response, while the all-particle matrix
$\widehat D_T^{\rm ind}\in\R^{N\times2N}$ supplies an independent physical
divergence diagnostic.

The row-restricted $D_T^q$ represents the terminal tangent in the selected
pressure-test space.  Consequently, $D_T^qu_T=0$ is an exact algebraic
statement in $Q_{h,0}^*$, while
\eqref{eq:discrete_terminal_tangent_consistency}, the all-particle independent
divergence, and the two area measures determine its spatial fidelity to the
continuum tangent and state equations.

For a periodic domain the corresponding pressure space is the mass-weighted
mean-zero quotient.  If $Z_p$ spans
\begin{equation}
  Q_{h,0}=\{p_h:\bm 1^TM_Qp_h=0\},
  \label{eq:periodic_mean_zero_space}
\end{equation}
then the response is formed on the quotient.  This construction retains every
divergence row and removes only the constant pressure gauge.

The free-boundary row restriction defines the concrete pressure-test space.
Rank, scaled singular values, manufactured $G$, $D$, and $DG$ actions, and
refinement establish its admissibility on every reported cloud.  Directly
testing the composition is essential because separate polynomial reproduction
of $D_h$ and $G_h$ leaves the cross-response accuracy undetermined.

\subsection{Discrete operator consistency}
\label{sec:discrete_operator_consistency}

Let $I_h$ denote particle sampling, let $I_h^Q$ map smooth pressures to their
coefficients, and let $\mathcal R_h^q$ restrict scalar fields to pressure-test
labels.  Consistency of the represented terminal tangent is measured by
\begin{equation}
 e_{D,h}(w)=
 \norm[M_Z]{D_{T,h}^q I_hw-\mathcal R_h^q(D_T^{\rm str}w)}
 \longrightarrow0
 \qquad (h\longrightarrow0)
 \label{eq:discrete_terminal_tangent_consistency}
\end{equation}
for smooth $w$ and a regular cloud family.  The composite pressure response
has the error decomposition
\begin{align}
 \varepsilon_{G,h}(q)
 &=G_{\Th,h}I_h^Qq-I_h(\mathcal G_\Th q),\notag\\
 \varepsilon_{DG,h}(q)
 &=D_{T,h}^qG_{\Th,h}I_h^Qq
   -\mathcal R_h^q(D_T^{\rm str}\mathcal G_\Th q)\notag\\
 &=\underbrace{D_{T,h}^qI_h(\mathcal G_\Th q)
   -\mathcal R_h^q(D_T^{\rm str}\mathcal G_\Th q)}_{\text{$D$-only error}}
   +\underbrace{D_{T,h}^q\varepsilon_{G,h}(q)}_{\text{propagated action error}}.
 \label{eq:discrete_composite_consistency}
\end{align}
The two limits are tested independently.  In particular, an $O(h^{-1})$
bound for the discrete derivative can convert a second-order action error into
a first-order propagated response error.  Case~3 reports both contributions,
including their boundary-layer decomposition.

The pressure solve enforces $D_{T,h}^qu_h^T=0$ in the selected test space.
The quantities $J_h-1$, $A_J$, and $A_\Gamma$ measure the associated
finite-step state geometry.  Together, these diagnostics separate represented
tangent closure, continuum consistency, and volume-state accuracy.

\subsection{Relative transport of the accepted operators}

At the beginning of step $n$, the accepted discrete state is
\begin{equation}
  \mathcal S_n=(x_n,u_n,D_n^q,G_n,\Grad_{{\rm geom},n}).
  \label{eq:production_state}
\end{equation}
Given the relative matrices $H_T,H_\Th$ in
\eqref{eq:relative_terminal_F}--\eqref{eq:relative_action_F}, their blockwise
Piola and cofactor transports are
\begin{align}
 D_T^qv
 &=j_T^{-1}D_n^q\bigl(j_TH_T^{-1}v\bigr),
 \label{eq:discrete_divergence_transport}\\
 G_\Th p
 &=j_\Th H_\Th^{-T}G_np,
 \label{eq:discrete_action_transport}\\
 G_Tp
 &=j_TH_T^{-T}G_np.
 \label{eq:discrete_terminal_action_transport}
\end{align}
The matrices $jH^{-1}$ and $jH^{-T}$ act particlewise on the velocity
blocks.  Formula \eqref{eq:discrete_action_transport} treats $G_n$ as the
complete current pressure-action state.  The recursive update therefore
applies the relative cofactor factor exactly once.

The geometry derivative reconstructs relative deformation and is transported
by the vector-gradient
chain rule,
\begin{equation}
 \Grad_{{\rm geom},n+1}v
 =\Grad_{{\rm geom},n}v\,H_T^{-1}.
 \label{eq:geometry_gradient_transport}
\end{equation}
The accepted divergence is committed directly through
\eqref{eq:full_operator_commit}, preserving the terminal constraint across
steps.

\subsection{Riesz scaling and the frozen pressure solve}

Let $m_i^0>0$ be the initial material quadrature weights.  At the beginning of
step $n$, the represented volume is
$m_i^n=m_i^0J_{{\rm cum},i}^n$ and its terminal candidate is
$m_i^T=m_i^nj_{T,i}$.  On a free-boundary pressure/test label
$a\in\mathcal I_q$, the Riesz maps used in every reported spectrum are
\begin{equation}
 M_Q=\diag(m_a^0)_{a\in\mathcal I_q},
 \qquad
 M_Z=\diag(m_a^T)_{a\in\mathcal I_q}.
 \label{eq:free_boundary_riesz_maps}
\end{equation}
Thus $M_Q$ measures the fixed material pressure coefficients and $M_Z$ measures
returned terminal constraint values; they are constructed separately even when
their entries happen to agree.  On a periodic quotient they are replaced by
$Z_p^T\diag(m_i^0)Z_p$ and $Z_D^T\diag(m_i^T)Z_D$, respectively.  Here $Z_p$
is the mass-mean-zero pressure basis and $Z_D$ removes only the verified global
compatibility direction.  No empirically fitted diagonal scaling is used.

The velocity Riesz map is
\begin{equation}
 M_U^{(n)}=\diag(\mu_1^n,\mu_1^n,\ldots,\mu_N^n,\mu_N^n),
 \label{eq:velocity_mass_matrix}
\end{equation}
and is held fixed during a pressure correction and its kinetic-work identity.
The affine and differential-vortex calculations use the fixed material
weights $\mu_i^n=m_i^0$.  The diagnostic choice does not enter the square
strong pressure equation.

The constraint matrix $D_T^q$ returns point values, so its scaled response is
\begin{equation}
 \widehat L_{T,\Th,h}
 =M_Z^{1/2}D_T^qG_\Th M_Q^{-1/2}.
 \label{eq:riesz_scaled_response}
\end{equation}
If dual constraint coefficients $D_{T,{\rm dual}}^q=M_ZD_T^q$ are stored,
the equivalent expression is
$M_Z^{-1/2}D_{T,{\rm dual}}^qG_\Th M_Q^{-1/2}$.
We report
\begin{equation}
 \beta_{T,\Th,h}=\sigma_{\min}(\widehat L_{T,\Th,h}),
 \qquad
 \kappa_{\rm cross}=
 \frac{\sigma_{\max}(\widehat L_{T,\Th,h})}
      {\sigma_{\min}(\widehat L_{T,\Th,h})}.
 \label{eq:scaled_cross_metrics}
\end{equation}
The Riesz-scaled quantities in \eqref{eq:scaled_cross_metrics} are the
stability diagnostics used throughout.

The following elementary lemma records the finite-dimensional stability
transfer used by the square collocation solve.

\begin{lemma}[Discrete cross-response stability]
\label{prop:discrete_cross_response_stability}
Let $A_h=D_T^qG_\Th$ be square and let
$\beta_{T,\Th,h}>0$ be defined by \eqref{eq:riesz_scaled_response}--
\eqref{eq:scaled_cross_metrics}.  If $A_hp_h=b_h$, then for every comparison
coefficient $p_I$,
\begin{equation}
 \norm[M_Q]{p_h-p_I}
 \le \beta_{T,\Th,h}^{-1}
      \norm[M_Z]{b_h-A_hp_I}.
 \label{eq:discrete_pressure_stability_bound}
\end{equation}
Moreover,
\begin{equation}
 \norm[M_U]{G_\Th(p_h-p_I)}
 \le C_{G,h}\beta_{T,\Th,h}^{-1}
      \norm[M_Z]{b_h-A_hp_I},
 \quad
 C_{G,h}=\norm[2]{M_U^{1/2}G_\Th M_Q^{-1/2}}.
 \label{eq:discrete_action_stability_bound}
\end{equation}
\end{lemma}

\begin{proof}
Set $e=M_Q^{1/2}(p_h-p_I)$.  Then
$\widehat L_{T,\Th,h}e=M_Z^{1/2}(b_h-A_hp_I)$, and the first estimate follows
from the smallest singular value.  Applying the scaled action to $e$ gives the
second estimate.
\end{proof}

The gap controls coefficient error, but action error also involves a
derivative.  A sharper transfer constant is
\begin{equation}
 C_{{\rm resp},h}
 =\norm[2]{M_U^{1/2}G_\Th A_h^{-1}M_Z^{-1/2}},\qquad
 \norm[M_U]{G_\Th(p_h-p_I)}
 \le C_{{\rm resp},h}\norm[M_Z]{b_h-A_hp_I}.
 \label{eq:discrete_response_transfer}
\end{equation}
This follows directly from $p_h-p_I=A_h^{-1}(b_h-A_hp_I)$, and
$C_{{\rm resp},h}\le C_{G,h}/\beta_{T,\Th,h}$.  Uniform bounds on
$\beta^{-1}$ and $C_{{\rm resp},h}$, together with an $O(h^r)$ comparison
residual, would give $O(h^r)$ coefficient and action errors.  A measured
spectral gap alone does not supply the latter bound; in particular,
$C_{G,h}$ need not remain bounded under refinement.  The present spectral
tests establish invertibility on the tested configurations, while consistency
and full-solution errors are measured separately.

For analytic forcing $b_h=\mathcal R_h^q(D_T^{\rm str}\mathcal G_\Th q)$,
the errors in \eqref{eq:discrete_composite_consistency} give the exact identity
\begin{equation}
 G_{\Th,h}p_h-I_h(\mathcal G_\Th q)
 =\varepsilon_{G,h}(q)
  -G_{\Th,h}A_h^{-1}\varepsilon_{DG,h}(q).
 \label{eq:solved_action_error_decomposition}
\end{equation}
It is the inverse-weighted composite defect, rather than
$\varepsilon_{DG,h}$ alone, that enters the solved action.  Thus a lower
consistency rate for $DG$ neither proves nor excludes a higher rate for that
action.  The manufactured solves in Case~3 measure this transfer with analytic
right-hand sides.

For frozen geometry the pressure equation is
\begin{equation}
  D_T^qG_\Th p=\dt^{-1}D_T^qu^\star.
  \label{eq:discrete_pressure_equation}
\end{equation}
It is solved by sparse LU in Cases~1 and~3.  Case~2 uses the pressure-only
Newton--Krylov implementation specified below.  The projection identities in
\cref{prop:frozen_projection} are tested with the actual matrices and random
as well as manufactured vectors.

\subsection{Nonlinear geometry iteration and final consistency rebuild}

Starting from $x_T^{(0)}=x^\star$, one Picard iteration consists of
\begin{align}
 &H_T^{(k)},j_T^{(k)},H_\Th^{(k)},j_\Th^{(k)}
     \quad\hbox{from }x_T^{(k)},\nonumber\\
 &D_T^{q,(k)}G_\Th^{(k)}p^{(k+1)}
     =\dt^{-1}D_T^{q,(k)}u^\star,\label{eq:picard_linear_solve}\\
 &u^{(k+1)}=u^\star-\dt G_\Th^{(k)}p^{(k+1)},\nonumber\\
 &x_{\rm cand}^{(k+1)}
   =x^\star-\gamma\dt^2G_\Th^{(k)}p^{(k+1)}.
 \label{eq:picard_update_equations}
\end{align}
Relaxation is applied only to the position iterate.  Both normalized RMS and
maximum kinematic increments must converge.

After Picard convergence the geometry is rebuilt, the square pressure equation
is solved, the velocity and kinematic position are recomputed, and the
operators are rebuilt once more on that returned position.  Acceptance uses
the complete pressure-test residual, the independent kinematic residual, and
the raw determinants.  Typical double-precision gates are $10^{-10}$ in the
relative pressure residual, $2\times10^{-10}$ and $2\times10^{-9}$ in the RMS
and maximum represented divergence, and $10^{-9}$ in the normalized final
kinematic residual.  Every accepted determinant is evaluated without clipping
and satisfies $j>10^{-6}$.  Pseudo-inverses, diagonal shifts, filtered pressure
modes, deleted rows, and post-acceptance corrective projections are absent.

Only after this final consistency rebuild passes do we commit
\begin{equation}
 x_{n+1}=x_T,\qquad u_{n+1}=u_T,\qquad
 D_{n+1}^q=D_T^q,\qquad G_{n+1}=G_T,
 \label{eq:full_operator_commit}
\end{equation}
together with \eqref{eq:geometry_gradient_transport}.  Hence
\begin{equation}
 D_{n+1}^qu_{n+1}=D_T^qu_T
 \label{eq:start_step_continuity}
\end{equation}
before any correction in the next step.  Recursive operators are also compared
with their direct cumulative Piola/cofactor transports.  These comparisons are
verification diagnostics; the recursive accepted matrices remain the
time-stepping state.

\begin{ccopalgorithm}[One accepted CCOP step]
\label{alg:accepted_ccop_step}
The complete step is listed in the order in which its configuration dependence
enters.
The nonpressure predictor is a specified part of each problem rather than an
implicit solver option.  Case~1 is force free and uses
\begin{equation}
 u^\star=u_n,\qquad x^\star=x_n+\dt u_n.
 \label{eq:force_free_predictor}
\end{equation}
For $f(x)=-A_fx$, with $A_f=\Omega^2I$ in Case~2 and $A_f=I$ in Case~3,
the implicit-midpoint force predictor is
\begin{equation}
 \left(I+\frac{\dt^2}{4}A_f\right)u_{n+1/2}
 =u_n-\frac{\dt}{2}A_fx_n,\qquad
 u^\star=2u_{n+1/2}-u_n,\qquad
 x^\star=x_n+\frac{\dt}{2}(u_n+u^\star).
 \label{eq:conservative_force_predictor}
\end{equation}
Starting from the incoming state \eqref{eq:production_state}, one step is:
\begin{enumerate}[label=\textbf{\arabic*.},leftmargin=2.2em,itemsep=0.25em]
\item Evaluate the problem-specific predictor and initialize $x_T^{(0)}$.
\item At Picard iterate $x_T^{(k)}$, reconstruct $H_T^{(k)}$ and the linear
path $H_\Th^{(k)}$ with $\Grad_{{\rm geom},n}$; evaluate raw determinants and
reject the trial if either is nonpositive or below the admissibility threshold.
\item Transport the incoming accepted operators by
\eqref{eq:discrete_divergence_transport}--
\eqref{eq:discrete_terminal_action_transport}.
\item Solve the square pressure-test equation
\eqref{eq:discrete_pressure_equation}, and evaluate
$u^{(k+1)}$ and $x_{\rm cand}^{(k+1)}$ from
\eqref{eq:picard_update_equations}.  Relax only the position iterate and repeat
until both RMS and maximum geometry increments converge.
\item Perform the final consistency rebuild: reconstruct the terminal and action geometry
from the returned position, resolve pressure, recompute velocity and position,
and rebuild once more on that accepted candidate.  Check the complete square
pressure residual, the represented terminal pressure-test residual, the
kinematic fixed point, and every raw Jacobian.
\item Commit \eqref{eq:full_operator_commit} and
\eqref{eq:geometry_gradient_transport} directly.  The pressure labels remain
fixed, and the accepted $D_T^q$ becomes the next-step divergence without a
post-step projection or scalar-gradient reconstruction.
\end{enumerate}
\end{ccopalgorithm}

The algebraic structure per nonlinear iteration is independent of $\Th$:
every placement reconstructs one terminal/action geometry pair and solves one
square pressure system of the same dimension.  Midpoint placement therefore
obtains the cancellation in \eqref{eq:placed_impulse_defect} without a second
pressure stage or an additional pressure equation.  Total nonlinear work can
still vary through the Picard count and is therefore recorded rather than
assumed equal.  In the affine multistep matrix, the 540 accepted steps for each
placement required mean Picard counts 4.213 for $\Th=1/2$ and 4.220 for
$\Th=1$ (both ranged from 3 to 6), so the realized nonlinear work is also
essentially identical in this mechanism test.

\subsection{Independent physical diagnostics}

The represented constraint, the independent divergence, and the area measures
are kept distinct.  The degree-three operator provides
\begin{equation}
 \widehat D_T^{\rm ind}u_T,
 \label{eq:independent_divergence}
\end{equation}
reported in the bulk, first boundary layer, boundary, and globally.  It never
appears in \eqref{eq:discrete_pressure_equation}.  The Jacobian area and the
boundary area are
\begin{equation}
 A_J=\sum_i m_iJ_{{\rm cum},i},
 \qquad
 A_\Gamma=\frac12\oint_\Gamma(x\,dy-y\,dx).
 \label{eq:two_area_measures}
\end{equation}
Their reported relative defects are
\begin{equation}
 e_A^J(t)=\frac{|A_J(t)-A_J(0)|}{A_J(0)},
 \qquad
 e_A^\Gamma(t)=\frac{|A_\Gamma(t)-A_\Gamma(0)|}{A_\Gamma(0)}.
 \label{eq:relative_area_defects}
\end{equation}
For the supplementary transport spot check we also use
\begin{equation}
 \eta_F(t)=
 \left(\frac{\sum_i m_i\norm[\mathrm F]{F_{{\rm cum},i}(t)-I}^2}
 {\sum_i m_i}\right)^{1/2},
 \qquad
 e_J^{\rm rms}(t)=
 \left(\frac{\sum_i m_i(J_{{\rm cum},i}(t)-1)^2}
 {\sum_i m_i}\right)^{1/2}.
 \label{eq:cumulative_deformation_diagnostics}
\end{equation}
Finally, $D_T^qu_T=0$ is the represented terminal pressure-test tangency
condition motivated by \cref{prop:state_tangent_identity}.  It is not relabeled
as a discrete $J_h=1$ constraint.  The two area measures in
\eqref{eq:two_area_measures} independently quantify the finite-volume state
consequence, including both spatial and temporal state errors rather than
only the algebraic closure tolerance.

\subsection{Constant-field and translation-equivariance diagnostics}
\label{sec:constant_field_diagnostics}

Piola transport of a consistent continuous divergence annihilates every
constant physical velocity.  The matrix recursion
\eqref{eq:discrete_divergence_transport} does not inherit this property solely
from recursive/direct agreement, because a discrete GMLS derivative need not
satisfy a product rule or commuting mixed derivatives.  We therefore monitor
the necessary constant-field defect
using
$\norm[M_Z,{\rm rms}]{z}^2=z^TM_Zz/(\bm1^TM_Z\bm1)$ and
$\norm[M_U,{\rm rms}]{v}^2=v^TM_Uv/\sum_i\mu_i$:
\begin{equation}
 e_{\rm const}(T)
 =\max_{k=1,2}\norm[M_Z,{\rm rms}]{D_T^qe_k},
 \qquad e_1=(1,0)^T,\quad e_2=(0,1)^T,
 \label{eq:constant_field_divergence_defect}
\end{equation}
with the constant vectors repeated at all particles.  For one directly
transported two-dimensional map reconstructed with $H_x,H_y$, expansion of
the discrete Piola rows gives
\begin{align}
 D_T^qe_1
 &=\mathcal R_h^q\!\left[J^{-1}(H_xH_y-H_yH_x)y\right],\notag\\
 D_T^qe_2
 &=-\mathcal R_h^q\!\left[J^{-1}(H_xH_y-H_yH_x)x\right].
 \label{eq:constant_field_commutator}
\end{align}
Thus exact annihilation is equivalent here to a discrete mixed-derivative
commutation property and is not implied by polynomial reproduction alone.

For the realized square pair, set
$\Pi_{\Th,h}^q=I-G_\Th(D_T^qG_\Th)^{-1}D_T^q$.  Its associated frozen
translation-equivariance defect is independent of the particular predictor
$v$:
\begin{equation}
 \begin{split}
 e_{\rm tr}(c)
 &=\norm[M_U,{\rm rms}]{\Pi_{\Th,h}^q(v+c)-\Pi_{\Th,h}^qv-c}\\
 &=\norm[M_U,{\rm rms}]{G_\Th(D_T^qG_\Th)^{-1}D_T^qc}.
 \end{split}
 \label{eq:frozen_translation_defect}
\end{equation}
These diagnostics test physical consistency of the transported rows; they are
separate from projector closure, which is exact for the matrix pair actually
solved.  Their nonaffine refinement is reported in
\cref{tab:supp_constant_translation_audit}.

\section{Numerical results}
\label{sec:numerical_results}

\subsection{Case 1: Terminal closure and pressure work are independent finite-step properties}
\label{sec:result1}

An affine solution lying in the polynomial reproducing space isolates the
mechanism without spatial-reconstruction contamination.  Following Roberts,
Shkoller, and Sideris
\cite{RobertsShkollerSideris2020}, let
$\B=\{X\in\R^2:|X|<1\}$ and
\begin{equation}
  x(t,X)=A(t)X,
  \qquad A(t)\in\mathrm{SL}(2,\R),
  \label{eq:affine_exact_map}
\end{equation}
where
\begin{equation}
 \ddot A=\lambda\cof A,
 \qquad
 \lambda=-\frac{2\det\dot A}{|A|_F^2},
 \qquad
 A(0)=I,\quad \dot A(0)=\diag(1,-1).
 \label{eq:affine_matrix_equation}
\end{equation}
The exact velocity and pressure are
\begin{equation}
 u(t,A(t)X)=\dot A(t)X,
 \qquad
 p(t,A(t)X)=\frac{\lambda(t)}2(1-|X|^2).
 \label{eq:affine_velocity_pressure}
\end{equation}
For these diagonal data, $A=\diag(a,a^{-1})$ and
\begin{equation}
 \lambda=\frac{2\dot a^2}{a^4+1},
 \qquad
 \ddot a=\frac{2\dot a^2}{a(a^4+1)},
 \qquad
 \dot a=\frac{\sqrt2a^2}{\sqrt{a^4+1}}.
 \label{eq:affine_scalar_reference}
\end{equation}
A DOP853 reference preserves $\det A=1$ to $1.11\times10^{-16}$ and the
matrix energy $|\dot A|_F^2/2=1$ to $6.66\times10^{-16}$.  The pressure is
positive in the disk and vanishes on its boundary.

The quasi-uniform disk levels have $h\simeq1/12,1/18,1/24,1/30$.  All four
square pressure responses are full rank; their complete scaled spectra and
pressure modes are reported in \cref{app:operator_audits}.  Because affine
trace-free velocity and quadratic pressure are reproduced, this problem
identifies finite-step structure rather than a spatial rate.

\subsubsection{The anchor determines which constraint is closed}

At $\Th=1/2$ we use the same predictor and placed action but form the pressure
equation with $D_n^q$, $D_\star^q$, or $D_T^q$.  For
$\dt=0.04,0.02,0.01,0.005$, each solve closes its own anchor.  The residual of
the same returned velocity under the terminal operator behaves as shown in
\cref{tab:anchor_rates}.  Current anchoring leaves an $O(\dt)$ terminal
residual and an $O(1)$ source mismatch.  Predictor anchoring is higher order in
this particular problem, whereas terminal anchoring reaches the
floating-point solve floor.

\begin{table}[t]
\centering
\caption{Observed rates in the affine anchor experiment.  The source mismatch
is $\dt^{-1}\norm{D_T^qu_\alpha^T}$.  Rates at the floating-point closure floor
are not fitted.}
\label{tab:anchor_rates}
\begin{tabular}{lrrr}
\toprule
anchor & terminal residual & source mismatch & $J_T-1$ \\
\midrule
$D_n^q$       & 1.001 & 0.001 & 2.000 \\
$D_\star^q$   & 2.999 & 1.999 & 3.999 \\
$D_T^q$       & fp64 floor & fp64 floor & 3.998 \\
\bottomrule
\end{tabular}
\end{table}

The predictor rates $O(\dt^3)$ and $O(\dt^2)$ are specific to the symmetric
affine reference.  Their value is structural: even an accurately closed,
higher-order nonterminal anchor leaves a distinct terminal residual, whereas
$D_T$ anchoring closes the accepted target to the solver tolerance.

\begin{figure}[t]
\centering
\begin{minipage}[t]{0.49\linewidth}
\centering
\safegraphic{\figroot/test1_affine_free_boundary/figures/anchor_residuals.png}
\end{minipage}\hfill
\begin{minipage}[t]{0.49\linewidth}
\centering
\safegraphic{\figroot/test1_affine_free_boundary/figures/theta_energy_closure.png}
\end{minipage}
\par\medskip
\safegraphic[0.72\linewidth]{\figroot/test1_affine_free_boundary/figures/projector_geometry.png}
\caption{Mechanism identification in the affine free-boundary flow.  Top
left: a solution that closes a lagged anchor retains a terminal source.  Top
right: all placements close their respective represented terminal pressure-test
constraints to the same solver tolerance while their pressure-induced kinetic
changes differ.  Bottom:
changing the placed action changes both the mass-scaled distance from the
terminal projector and the principal angle between the placed and terminal
action ranges.  The third projector panel includes the auxiliary mass-adjoint
control, which imposes an orthogonal endpoint baseline by definition.}
\label{fig:test1_mechanism}
\end{figure}

\subsubsection{The placed action determines projector geometry and pressure work}

With the terminal target fixed, a sweep over
$\Th=0,0.05,\ldots,1$ at $\dt=0.02$ shows that every placement satisfies
$D_T^qu_T=0$ to the common solver tolerance, and the independently evaluated
work--energy identity \eqref{eq:discrete_pressure_work_energy} closes at the
floating-point arithmetic floor.  Thus every placement is exactly
work-consistent with its own represented action, although those actions
produce different finite-step work.

A separate frozen projector experiment holds the terminal geometry returned by the
$h\simeq1/12$, $\Th=1/2$ case fixed and sweeps the action configuration along
that accepted position path.  It tests exactly the operator family
in \eqref{eq:frozen_ccop_operator} and separates projector geometry from the
change of nonlinear terminal state across independently solved placements.  Across
all 21 frozen placements, closure, idempotence, and action annihilation hold
to the direct-solve tolerance; their numerical floors are listed in the
appendix.  Relative to the terminal projector, the
mass-scaled projector distance and largest action-space angle are
\begin{equation}
\begin{array}{c|ccc}
\Th & 0 & 1/2 & 1 \\ \hline
\norm[M_U]{\Pi_\Th-\Pi_T}
 & 3.1319\times10^{-2} & 1.5669\times10^{-2} & 1.68\times10^{-14}\\
\angle_{\max}(\range G_\Th,\range G_T)
 & 1.14565^\circ & 0.57285^\circ & 3.82\times10^{-6}\,{}^\circ.
\end{array}
\label{eq:test1_projector_geometry_values}
\end{equation}
The measured endpoint discretization baseline is
$\norm[M_U]{\Pi_T}=1.61863$, with
$e_{{\rm Gr},{\rm full}}=0.2484$ under the explicit normalization
\eqref{eq:full_space_green_compatibility}.
This value measures the full-space mismatch between the independently
constructed strong pressure-test divergence and the mass adjoint induced by
$G_T$.  Because it compares two discrete realizations, it quantifies a
compatibility baseline; separate manufactured tests assess approximation
accuracy.  The projector identities require
invertibility of $D_T^qG_T$, not discrete mass adjointness, and remain at the
direct-solve floor.
The distance and angle in \eqref{eq:test1_projector_geometry_values} isolate
the increment produced by configuration placement.  The baseline is fixed as $\Th$ varies,
whereas the added angle at $\Th=0$ is $0.019995$ radians at $\dt=0.02$, the
$O(\dt)$ scale predicted by \cref{prop:projector_proximity}.  These diagnostics
therefore measure different effects rather than competing percentages.

The auxiliary control \eqref{eq:compatible_endpoint_control} makes the
identification independent of that baseline.  It gives
\begin{equation}
 \begin{aligned}
 \norm[M_U]{\Pi_T^{\rm ad}}&=1.0000000\quad\hbox{(to roundoff)},&
 \norm[M_U]{\Pi_0^{\rm ad}}&=1.00019994,\\
 \norm[M_U]{\Pi_0^{\rm ad}-\Pi_T^{\rm ad}}
 &=1.9998\times10^{-2}.&&
 \end{aligned}
 \label{eq:test1_compatible_projector_control}
\end{equation}
with endpoint mass-self-adjoint defect $1.38\times10^{-15}$.  Closure,
idempotence, and action annihilation remain at $2.3\times10^{-15}$ or below.
Thus a compatible endpoint is orthogonal to roundoff, whereas moving the same
action space away from $T$ produces a nonzero projector distance and a norm
strictly larger than one.  The meshless sweep and the compatible control
therefore identify, respectively, the realized total projector and the pure
configuration-induced increment.

The full-space mismatch is not active on the polynomial pair excited by this
affine solution.  Taking $\mathcal V_r$ as the trace-free affine Eulerian
velocity space and
$\mathcal Q_r=\operatorname{span}\{1-|X|^2\}$ gives
$e_{{\rm Gr},r}=2.26\times10^{-16}$ in
\eqref{eq:restricted_green_compatibility}; the mapped pressure-action error is
also below $10^{-15}$.  The nearly antisymmetric endpoint work values are
consistent with this restricted compatibility, rather than with small
full-space nonorthogonality.

The frozen common target therefore coexists with genuinely different
correction ranges.  The nonlinear placement sweep likewise exhibits common
represented closure and different work.  This is qualitatively consistent with
the sensitivity identity \eqref{eq:frozen_work_placement_sensitivity};
projector distance, action-space angle, and work provide complementary views
of the same range variation.  The refined endpoint work values approach
\begin{equation}
 \Delta K_p(0)=+3.142\times10^{-4},
 \qquad
 \Delta K_p(1)=-3.140\times10^{-4},
 \label{eq:test1_endpoint_work_values}
\end{equation}
while $\Delta K_p(1/2)\simeq-6.28\times10^{-8}$.  The observed sign bracket
admits a work-neutral root,
\begin{equation}
 \Th^\star=0.4999001,
 \qquad
 |\Th^\star-1/2|=O(\dt^{1.999}).
 \label{eq:test1_work_neutral_placement}
\end{equation}
The endpoint sign change brackets this root.  Work neutrality has a physical reference here: the
exact force-free affine solution has constant kinetic energy, so its net
pressure work is zero.  The convergence of $\Th^\star$ to $1/2$ is therefore a
pressure-impulse consistency test, not merely a root-finding diagnostic.

The stage hypothesis in \cref{prop:placed_impulse_order} is tested separately
from accumulated trajectory error.  Starting from the exact affine state at
$t_n=0.2$, we solve one coupled step for
$\dt=0.02,0.01,0.005,0.0025$ and compare the returned action with
$g(t_n+\Th\dt)$.  The rates in \cref{tab:affine_stage_rates} show the needed
second-order midpoint stage consistency and the generic first-order endpoint
behavior.  Applying the exact quadratic pressure through the numerical
action-path operator gives a maximum RMS error
$8.5\times10^{-16}$, so the rates measure temporal stage placement rather
than spatial action reproduction.  The action-position errors themselves have
rates 1.992 at $\Th=1/2$ and 2.994 at $\Th=1$; the first-order endpoint action
error is therefore traced to the pressure coefficient stage, not to the
geometric path or the action reconstruction.

\begin{table}[t]
\centering
\caption{Exact-start placed-action consistency at $t_n=0.2$.  Here
$e_{g,\Th}=\norm[M_U]{g_\Th^h-g(t_n+\Th\dt)}$ and $e_{p,\Th}$ is the
corresponding pressure-coefficient error.}
\label{tab:affine_stage_rates}
\begin{tabular}{lrr}
\toprule
$\Th$ & rate of $e_{g,\Th}$ & rate of $e_{p,\Th}$ \\
\midrule
$0$   & 1.001 & 1.001 \\
$1/2$ & 1.967 & 1.969 \\
$1$   & 1.010 & 1.007 \\
\bottomrule
\end{tabular}
\end{table}

Multistep advancement to $T=0.5$ turns the one-step work separation into the
cumulative consequence reported in \cref{tab:affine_time_rates}.  Midpoint
placement is approximately second order for
this affine solution, whereas terminal placement is approximately first order
in the principal state and energy errors.  These rates validate
the fixed-space mechanism in
\cref{prop:actual_ccop_local_defect}: the force-free stage and trapezoidal
kinematics cancel the leading configuration-induced velocity defect at
$\Th=1/2$, while $\Th=1$ retains it.  The separate placed-action rates in
\cref{tab:affine_stage_rates} verify the hypothesis of
\cref{prop:placed_impulse_order} on this reproduced affine field.

\begin{table}[t]
\centering
\caption{Observed temporal rates for the affine multistep solution.
The last column is the pressure-coefficient error evaluated at the action
time, not the error of the placed vector action.}
\label{tab:affine_time_rates}
\begin{tabular}{lrrrrrr}
\toprule
$\Th$ & $x$ & $u$ & $F$ & $J$ & kinetic energy & $e_{p,\Th}$ \\
\midrule
$1/2$ & 2.000 & 2.000 & 2.000 & 2.000 & 2.000 & 1.985 \\
$1$   & 1.005 & 0.998 & 1.005 & 2.006 & 0.997 & 1.006 \\
\bottomrule
\end{tabular}
\end{table}

The affine result closes the terminal-target/action-placement mechanism established above.  Within each nonlinear placement,
$D_T^q[x_T(\Th)]$ selects the represented target on that placement's own
terminal state.  The frozen projector experiment then holds one such target fixed and isolates
how $G_\Th$ changes the correction range and hence the oblique projector used
to reach it.  The placed solves all attain their respective terminal
pressure-test targets but perform different finite-step pressure work.  The
multistep rates show that this one-step distinction accumulates into observable
state and energy differences.  Since the two placements have the same pressure
space and system dimension, require one pressure solve per Picard iteration,
and have essentially identical Picard counts in this test, the second-order
midpoint pressure path requires no additional pressure stage.


\subsection{Case 2: Long-time ISPH comparison under matched spatial discretization}
\label{sec:result2}

Case~1 separately identifies the terminal target and pressure-action placement
axes.  We now test the long-time consequence of replacing the classical
predictor-configuration ISPH projection by the complete CCOP architecture.
The two-dimensional conservative affine drop places the exact velocity and
pressure derivatives in the polynomial design spaces, thereby isolating the
long-time pressure update from spatial truncation.  The
initial circular drop has $R=0.5$, the forcing frequency is $\Omega=1.2$, and
the velocity field is affine,
\begin{equation}
 u_0=A_0x,\qquad v_0=-A_0y,
 \qquad A_0=0.4,
 \label{eq:case2_affine_initial_velocity}
\end{equation}
while the body force is $-\Omega^2x$.  The exact motion remains affine and the
pressure is quadratic, placing their spatial derivatives within the polynomial
design spaces of the operators below.  The small moment regularization in the
archived driver is specified in \cref{sec:case2_implementation_contract}.
Writing the principal semi-axes as
$a(t)$ and $b(t)=R^2/a(t)$, with $\delta=\dot a/a$, gives the scalar reference
\begin{equation}
 \dot a=\delta a,
 \qquad
 \dot\delta=
 \frac{R^4/a^4-1}{R^4/a^4+1}
 \bigl(\delta^2+\Omega^2\bigr),
 \qquad
 p_c=\frac{a^2}{2}
 \bigl(\dot\delta+\delta^2+\Omega^2\bigr).
 \label{eq:case2_affine_reference}
\end{equation}
An independent DOP853 integration of \eqref{eq:case2_affine_reference} supplies
the semi-axes, centre pressure, phase, and conservative energy exchange through
$T=30$.  This construction targets the accumulated pressure-update error.

\subsubsection{Matched ISPH--CCOP comparison}

We compare the same-discretization ISPH workflow $D_\star^qG_\star$ with CCOP
$D_T^qG_{1/2}$ at $R/\Delta x=50$.  The archived \texttt{SPH2}
driver uses a cubic SPH kernel to weight a quadratic MLS fit; the \texttt{MLS2}
and \texttt{SPH2} strict drivers provide two realizations of this shared
derivative design.  The pressure uses a general nodal field with homogeneous
Dirichlet boundary values.
The common fixed material derivatives realize
\begin{equation}
 D_\alpha^q=\mathcal R_h^q\!\left[
 J_\alpha^{-1}\Div_0^h(J_\alpha F_\alpha^{-1}\,\cdot\,)
 \right],\qquad
 G_\alpha^{\rm phys}p=\rho_0^{-1}F_\alpha^{-T}\Grad_0^hp.
 \label{eq:case2_physical_action}
\end{equation}
The second operator uses the native physical-pressure convention
$F_\alpha^{-T}\Grad_0^h$.  Both branches use this convention.  For a spatially
uniform affine Jacobian, the cofactor variant has the same range and frozen
projector, with a rescaled pressure coefficient.  Case~3 uses the full
cofactor convention for its nonaffine maps.

The baseline is a \emph{same-discretization ISPH (star--star) implementation}:
it follows the predictor--pressure--corrector workflow of projection-based
ISPH \cite{CumminsRudman1999,XuStansbyLaurence2009}, with both pressure
operators evaluated on the pressure-free predicted configuration.  It solves
\begin{equation}
 D_\star^qG_\star^{\rm phys}p_\star
  =\Delta t^{-1}D_\star^qu^\star,\qquad
 u_\star^T=u^\star-\Delta tG_\star^{\rm phys}p_\star,\qquad
 x_\star^T=x^\star-\tfrac12\Delta t^2G_\star^{\rm phys}p_\star.
 \label{eq:case2_same_discretization_isph}
\end{equation}
The particles, material stencils, pressure unknowns, boundary labels,
quadrature, force predictor, step size, and position-correction coefficient
are shared with CCOP.  The resulting matched comparison isolates the change
from the star--star ISPH workflow to terminal-target, midpoint-action CCOP.
The fixed star--star system is solved by sparse LU.  Its own-anchor residual
and the common terminal diagnostic $D_T^qu_\star^T$ are recorded before any
further correction.

CCOP couples the same terminal target to the midpoint action.  For a
trial pressure $p$, the corrected state is
\begin{equation}
 u_T(p)=u^\star-\Delta t\,G_{1/2}^{\rm phys}[x_T(p)]p,
 \qquad
 x_T(p)=x^\star-\frac{\Delta t^2}{2}\,G_{1/2}^{\rm phys}[x_T(p)]p,
 \label{eq:case2_terminal_update}
\end{equation}
and the action configuration is
\begin{equation}
 x_{1/2}(p)=\frac12\bigl(x^n+x_T(p)\bigr).
 \label{eq:case2_midpoint_action_configuration}
\end{equation}
The pressure is obtained from the nonlinear terminal residual
\begin{equation}
 \mathcal R(p)=D_T^q[x_T(p)]u_T(p)=0.
 \label{eq:case2_terminal_residual}
\end{equation}
The CCOP branch follows the pressure-dependent chain
\begin{equation}
 p\longmapsto x_{1/2}(p)
 \longmapsto G_{1/2}^{\rm phys}p
 \longmapsto (x_T(p),u_T(p))
 \longmapsto D_T^qu_T(p).
 \label{eq:case2_residual_chain}
\end{equation}
The implemented map is differentiated by a JAX JVP in the matrix-free
Newton--Krylov solve, including its unrolled geometry iterations.  Acceptance
separately checks the rebuilt terminal residual, raw positive Jacobians, and
the kinematic fixed-point defect.  Case~2 uses this pressure-only solver
variant; Cases~1 and~3 use the outer block-Picard realization.  Implementation
details are recorded in \cref{sec:case2_implementation_contract}.

\begin{figure}[t]
\centering
\safegraphic[0.98\linewidth]{\figroot/test2_affine_long_time/fig_case2_long_time_comparison.pdf}
\caption{\textbf{Long-time consequence of the complete CCOP architecture relative to matched ISPH.}
(a) Major (solid) and minor (dashed) semi-axes against the DOP853 affine
reference.  (b) Centre-pressure error evaluated at the action time owned by
each branch.  (c,d) Kinetic- and potential-energy errors, normalized by the
initial exact mechanical energy.  (e) Relative mechanical-energy drift.
(f) Terminal-divergence diagnostic evaluated on the returned state.  The
star--star ISPH and CCOP midpoint workflows use the same particles, pressure
space, material derivatives, force predictor, and position update.  ISPH uses
$\Delta t=2\times10^{-3}$; the three CCOP curves use
$\Delta t=4\times10^{-3},2\times10^{-3},10^{-3}$.}
\label{fig:case2_affine_transfer}
\end{figure}

Figure~\ref{fig:case2_affine_transfer} compares the classical matched ISPH
projection with the complete CCOP architecture.  At the common step size
$\Delta t=2\times10^{-3}$, the maximum mechanical-energy error decreases from
$4.31\times10^{-3}$ to $1.45\times10^{-7}$; the relative semi-axis errors
decrease from approximately $2.3\times10^{-3}$ to $4.6\times10^{-6}$, and the
centre-pressure error decreases from $1.93\times10^{-3}$ to
$4.13\times10^{-6}$.  The returned ISPH state retains a terminal diagnostic of
$5.36\times10^{-7}$ after closing its own predicted-configuration constraint,
whereas CCOP closes its terminal pressure-test constraint to the nonlinear
solve tolerance.  Case~1 identifies the separate target and action mechanisms;
the present comparison shows their combined long-time consequence.  Across the
three CCOP step sizes, the observed
least-squares rates are $2.000$ for both semi-axes, $1.994$ for the kinetic and
potential components, $1.988$ for the maximum mechanical-energy error, and
$2.028$ for the centre pressure.

\begin{table}[t]
\centering
\caption{Long-time affine diagnostics for the quadratic material-derivative
realization archived under the SPH2 driver label, at
$R/\Delta x=50$ and $T=30$.  Here $R_{\rm own}$ is the residual of the anchor
used by each branch, $R_T$ is the common terminal diagnostic, $e_a$ and $e_b$
are relative RMS errors of the major and minor semi-axes,
$E_{\rm err}=\max_t|E/E_0-1|$, and $e_{p_c}$ is the relative RMS centre-pressure
error at the action time, obtained by the stored local quadratic centre fit.
ISPH denotes the same-discretization star--star workflow of
\eqref{eq:case2_same_discretization_isph}.}
\label{tab:case2_transfer_numbers}
\scriptsize
\setlength{\tabcolsep}{3.5pt}
\renewcommand{\arraystretch}{1.35}
\begin{tabular}{lccccccc}
\toprule
Branch & $\Delta t$ & $R_{\rm own}$ & $R_T$ & $e_a$ & $e_b$ &
$E_{\rm err}$ & $e_{p_c}$ \\
\midrule
ISPH (star--star) & $2{\times}10^{-3}$ & $1.81{\times}10^{-15}$ & $5.36{\times}10^{-7}$ & $2.40{\times}10^{-3}$ & $2.27{\times}10^{-3}$ & $4.31{\times}10^{-3}$ & $1.93{\times}10^{-3}$ \\
CCOP midpoint & $4{\times}10^{-3}$ & $4.87{\times}10^{-13}$ & $4.87{\times}10^{-13}$ & $1.65{\times}10^{-5}$ & $1.83{\times}10^{-5}$ & $5.78{\times}10^{-7}$ & $1.67{\times}10^{-5}$ \\
CCOP midpoint & $2{\times}10^{-3}$ & $1.78{\times}10^{-14}$ & $1.78{\times}10^{-14}$ & $4.13{\times}10^{-6}$ & $4.57{\times}10^{-6}$ & $1.45{\times}10^{-7}$ & $4.13{\times}10^{-6}$ \\
CCOP midpoint & $1{\times}10^{-3}$ & $1.31{\times}10^{-14}$ & $1.31{\times}10^{-14}$ & $1.03{\times}10^{-6}$ & $1.14{\times}10^{-6}$ & $3.68{\times}10^{-8}$ & $1.00{\times}10^{-6}$ \\
\bottomrule
\end{tabular}
\end{table}

Table~\ref{tab:case2_transfer_numbers} gives the corresponding scalar
diagnostics.  Because the particles, pressure space, material stencils,
quadrature, force predictor, and kinematic update are matched, the comparison
isolates the long-time consequence of changing from the classical star--star
projection workflow to the terminal-target/midpoint-action CCOP architecture.
The one-step attribution of those two configuration roles remains the task of
Case~1.  Equal pressure dimension does not imply equal cost for the frozen LU
and nonlinear solves.

\subsubsection{Relation to other particle formulations}

Particle regularization, continuity corrections, and consistent SPH gradients
address complementary error sources
\cite{Wang2019IPST,GotohKhayyerGotoh2024,ZhangHighOrderSPH}.
Published central-force drop histories are not a matched performance test
for the present run: for example, the RKGC case in
\cite[section~5.4]{ZhangHighOrderSPH} uses $A_0/\Omega=1$, whereas here
$A_0/\Omega=1/3$.  Normalizing energy by $E_0$ does not remove that difference
in deformation amplitude.  The algorithmic evidence in this case is therefore
the same-discretization ISPH--CCOP comparison in
\cref{tab:case2_transfer_numbers}, rather than a ranking of unmatched
published energy curves.

\subsection{Case 3: Configuration separation persists under nonaffine material shear}
\label{sec:result3}

The differential vortex establishes the same construction under a spatially
varying material map.  On the unit disk, set
\begin{equation}
 \Omega(r)=1-r^2,
 \qquad
 \Phi_t(X)=\mathcal R\!\left(t\Omega(|X|)\right)X,
 \qquad
 u(x)=\Omega(|x|)Kx,
 \quad
 K=\begin{pmatrix}0&-1\\1&0\end{pmatrix}.
 \label{eq:differential_vortex_map}
\end{equation}
The boundary is stationary because $\Omega(1)=0$, but different radii rotate
through different angles.  With $a=t\Omega(|X|)$,
\begin{equation}
 F_{\rm ex}(X,t)=\mathcal R(a)
 \left[I+(KX)\otimes\nabla_Xa\right],
 \qquad
 \det F_{\rm ex}=1.
 \label{eq:differential_vortex_F}
\end{equation}
The determinant follows from the matrix determinant lemma and
$\nabla_Xa\cdot KX=0$.  In a polar frame the nonrotational factor is a shear
with $s=-2tr^2$, giving
\begin{equation}
 \kappa_2(F_{\rm ex})
 =\left(\frac{\sqrt{s^2+4}+|s|}{2}\right)^2\le5.83
 \quad\text{at }T=1.
 \label{eq:differential_vortex_condition}
\end{equation}
With the conservative body force $f=-x$, radial momentum balance yields
\begin{equation}
  p(r)=\frac12(1-r^2)-\frac16(1-r^2)^3,
  \label{eq:differential_vortex_pressure}
\end{equation}
and direct differentiation verifies
$(u\cdot\nabla)u=-\Omega(r)^2x=-\nabla p+f$.

\subsubsection{Boundary-action propagation limits composite consistency}

Prescribing the exact particle positions isolates spatial reconstruction from
time integration and nonlinear iteration.  The analytic
reference is evaluated on the same linear action path
\eqref{eq:linear_action_path} used by the algorithm.  At shear times
$\chi=0.5,1,1.5$, placements
$\Th=0,1/2,1$, and four spatial levels
$h\simeq1/12,1/18,1/24,1/30$, we measure
\begin{equation}
  e_F,\qquad e_J,\qquad e_G,\qquad e_D,\qquad e_{DG}.
  \label{eq:test2_principal_errors}
\end{equation}
The action tests include the exact physical pressure and nonsymmetric
Dirichlet pressures.  The divergence tests include the vortex velocity and an
independent trace-free field.  The composite $DG$ reference is evaluated
directly, so the central cross response receives its own convergence test.

At the principal shear $\chi=1$ and $\Th=1/2$, the rates in
\cref{tab:test2_spatial_rates} show approximately second-order geometry and
separate operator actions, while the composite response exposes a
boundary-localized contribution with a lower observed order.  Every
square response remains full rank, and the minimum Riesz-scaled cross gap over
the entire prescribed-map study is $0.4615$.

\begin{table}[t]
\centering
\caption{Global mass-weighted prescribed-map rates at $\chi=1$ and
$\Th=1/2$.}
\label{tab:test2_spatial_rates}
\begin{tabular}{lrr}
\toprule
quantity & fitted rate & finest error \\
\midrule
$e_F$ & 2.013 & $4.045\times10^{-3}$ \\
$e_J$ & 2.281 & $2.753\times10^{-3}$ \\
$e_G$ & 1.999 & $1.118\times10^{-3}$ \\
$e_D^q$ & 1.869 & $1.910\times10^{-3}$ \\
$e_D^{\rm ind}$ & 3.128 & $3.582\times10^{-5}$ \\
$e_{DG}^q$ & 1.153 & $2.056\times10^{-2}$ \\
$e_{DG}^{\rm ind}$ & 1.266 & $4.825\times10^{-2}$ \\
\bottomrule
\end{tabular}
\end{table}

\begin{figure}[t]
\centering
\begin{minipage}[t]{0.49\linewidth}
\centering
\safegraphic{\figroot/test2_differential_vortex/figures/part2a_principal_convergence.png}
\end{minipage}\hfill
\begin{minipage}[t]{0.49\linewidth}
\centering
\safegraphic{\figroot/test2_differential_vortex/figures/part2c_radial_phase.png}
\end{minipage}
\caption{Nonaffine extension of the configuration split.  Left: prescribed-map
convergence of $F,J,G,D$, and $DG$.  Right: exact and computed angular
displacement as a function of radius for the complete time-stepping evolution.}
\label{fig:test2_main}
\end{figure}

The composite rate localizes the leading spatial contribution.  Decomposing
the independent cross error into a $D$-only term and the error propagated from
$G$ shows that placed-action propagation dominates near the boundary.  The
first-layer propagated-$G$ sequence is
\begin{equation}
  0.05825,\quad0.05531,\quad0.04809,\quad0.04337,
  \label{eq:test2_boundary_G_sequence}
\end{equation}
Its four-level fitted rate is $0.326$, and an additional $h=1/36$ value is
$0.03824$.  The corresponding first-layer share of the interior
pressure-test labels decreases as $0.165,0.111,0.0835,0.0662$, consistent with
an $O(h)$-thick boundary region.  For a regional RMS error of
order $h^r$, an $O(h)$ mass fraction contributes only a factor $h^{1/2}$
to its global RMS norm.  The inverse response in
\eqref{eq:discrete_response_transfer} further transforms this localized
operator error.  The data therefore identify boundary propagation of the
placed action as the leading limitation on composite consistency, while the
smooth-vortex solution retains its separate convergence rate.

\subsubsection{The inverse response retains accurate smooth pressure actions}

To measure what this consistency defect does to the solved pressure, we
freeze the same prescribed configurations at $\chi=1$, $\Th=1/2$ and solve
\begin{equation}
 A_hp_h=b_{\rm ex},\qquad
 b_{\rm ex}=\mathcal R_h^q(D_T^{\rm str}\mathcal G_{1/2}q).
 \label{eq:manufactured_composite_solve}
\end{equation}
The analytic composite action supplies the right-hand side, producing a
nontrivial manufactured pressure solve.  All particles, stencils, pressure
spaces, and operators are held fixed.  For the physical
radial pressure, \cref{tab:composite_pressure_transfer} shows that the
near-first-order composite residual gives pressure and final action errors
with fitted rates 2.183 and 1.990.  The directly measured transfer constant
$C_{{\rm resp},h}$ decreases over these four clouds.  The total action error is
the sum in \eqref{eq:solved_action_error_decomposition}, which includes direct
action error and the inverse-transmitted composite defect.  For these data,
the inverse-transmitted $D$-only and propagated-action contributions converge
at rates 1.978 and 1.991, respectively.  The improvement therefore occurs before
the two terms are combined.  Their separate convergence rules out final-error
cancellation and boundary-region contraction as the sole explanation.

\begin{table}[t]
\centering
\caption{Frozen manufactured pressure solve at $\chi=1$, $\Th=1/2$.
Errors are mass-weighted RMS values for the physical pressure.  The action
error compares $G_{1/2,h}p_h$ with the analytic placed action.
$C_{{\rm resp},h}$ uses the Riesz scaling in
\eqref{eq:discrete_response_transfer}.}
\label{tab:composite_pressure_transfer}
\begin{tabular}{@{}rrrrr@{}}
\toprule
$1/h$ & $e_{DG}$ & $e_{p,\rm solve}$ & $e_{G p,\rm solve}$ & $C_{{\rm resp},h}$ \\
\midrule
12 & $5.925\times10^{-2}$ & $3.487\times10^{-3}$ & $6.867\times10^{-3}$ & 1.154 \\
18 & $3.826\times10^{-2}$ & $1.408\times10^{-3}$ & $3.052\times10^{-3}$ & 0.964 \\
24 & $2.718\times10^{-2}$ & $7.570\times10^{-4}$ & $1.725\times10^{-3}$ & 0.871 \\
30 & $2.056\times10^{-2}$ & $4.727\times10^{-4}$ & $1.108\times10^{-3}$ & 0.822 \\
\midrule
Fitted rate & 1.153 & 2.183 & 1.990 & --- \\
\bottomrule
\end{tabular}
\end{table}

The nonsymmetric pressures $q_1=(1-|X|^2)X_1$ and
$q_2=(1-|X|^2)X_1X_2$ give solved-pressure rates 2.003 and 2.000,
and final-action rates 1.988 and 1.961.  All twelve solves have relative
linear residual below $3.9\times10^{-14}$.  These results support accurate
response for the tested smooth fields and configurations.  Arbitrary pressure
data and untested shear ranges remain outside this measured conclusion.

The stronger prescribed shear also retains an invertible composite response.  At
$\chi=1.5$, where $\max\kappa_2(F_{\rm ex})=10.91$, the midpoint responses
remain full rank; from $h\simeq1/12$ to $1/30$, the Riesz-scaled gap changes
from $0.653$ to $1.246$ while $e_{DG}^q$ decreases from $8.17\times10^{-2}$ to
$2.73\times10^{-2}$.  Thus the cross response remains stably invertible and
convergent up to the tested condition number.  The full table is given in
\cref{tab:supp_strong_shear}.

\paragraph{Transport preserves the reconstruction accuracy.}

An exact-increment verification divides $[0,1]$ into
$4,8,16,32,64$ analytic subintervals.  Relative deformation is reconstructed
from exact consecutive configurations, while $D$ and $G$ are transported by
the same matrix recursions used in time stepping.  Recursive and direct
cumulative operators agree at
$4.7\times10^{-16}$--$1.43\times10^{-15}$, far below the
$10^{-2}$--$10^{-3}$ direct spatial reconstruction error.  The complete depth
study is given in Supplementary Fig.~\ref{fig:supp_transport}.  Recursive
transport therefore preserves the spatial accuracy established by the direct
nonaffine reconstruction.

\subsubsection{Complete nonaffine convergence}

The complete method advances to $T=1$ with its conservative-force
predictor, nonlinear terminal geometry, placed pressure action, square
terminal constraint, and full-operator commit.  At fixed $\dt=1/1024$, both
principal placements converge at approximately second order in the state,
deformation, radial phase, and independent divergence; selected rates are
reported in \cref{tab:test2_production_rates}.

\begin{table}[t]
\centering
\caption{Spatial rates for the complete differential-vortex evolution.}
\label{tab:test2_production_rates}
\begin{tabular}{lrrrrrrr}
\toprule
$\Th$ & $x$ & $u$ & $p$ & $F$ & $J$ & phase & independent $D$ \\
\midrule
$1/2$ & 1.994 & 1.998 & 2.001 & 1.927 & 2.124 & 1.995 & 2.034 \\
$1$   & 2.067 & 2.074 & 2.062 & 1.936 & 2.124 & 2.142 & 2.034 \\
\bottomrule
\end{tabular}
\end{table}

\begin{figure}[t]
\centering
\safegraphic[0.98\linewidth]{\figroot/test2_differential_vortex/figures/fig_case3_nonaffine_state_energy_divergence.pdf}
\caption{\textbf{State, energy, and closure under nonaffine material shear.}
(a) Radial angular displacement at $T=1$ for the finest cloud.  (b) Final
spatial errors at $\Delta t=1/1024$; circles, squares, triangles, and diamonds
denote position, velocity, pressure, and independent-divergence errors,
respectively.  Solid blue and dashed orange curves denote midpoint and terminal
placement.  (c--e) Relative kinetic-, conservative-potential-, and
mechanical-energy invariant errors on the $R/h=30$ cloud.  (f) Represented
terminal closure (solid) and the independent strong-divergence diagnostic
(dotted) on the same cloud.}
\label{fig:case3_state_energy_divergence}
\end{figure}

Both placements achieve approximately second-order state accuracy under
strictly nonaffine material shear.  Together with the prescribed-map and
exact-increment results, this establishes that the target/action separation,
cross response, and committed operators survive spatially varying material
deformation.  Figure~\ref{fig:case3_state_energy_divergence} shows the same
conclusion dynamically: the radial phase follows the analytic differential
rotation, the physical state and independent divergence decrease under spatial
refinement, and represented terminal closure remains at the algebraic floor.
Because this exact map preserves the circular boundary, major and minor axes
are identically non-discriminating; radial phase is the deformation observable.

The nonaffine case also supplies a fixed-cloud temporal-order check.  At
$R/h=24$, the exact-error curves for midpoint placement are nearly unchanged
over $\dt=1/64,1/128,1/256,1/512$.  Using common material labels, successive
differences $\norm{U_{\dt}-U_{\dt/2}}$ measure approach to a common
fixed-cloud limit while cancelling its time-independent spatial bias.
The resulting
mass-weighted rates are
\begin{equation}
\begin{array}{c|ccccc}
\Th & x & u & F & J & \hbox{radial phase}\\ \hline
1/2 & 2.000 & 2.000 & 2.000 & 1.993 & 1.998\\
1   & 0.999 & 1.001 & 1.001 & 1.003 & 0.997.
\end{array}
\label{eq:test2_temporal_self_difference_rates}
\end{equation}
The midpoint pressure-coefficient self-difference has rate 1.538, while the
terminal coefficient has rate 1.015.  The stored midpoint differences on
three clouds are consistent with a mixed $a_h\dt+b_h\dt^2$ norm model:
$a_h$ decreases at rate 2.007 while $b_h$ changes by less than $0.25\%$
(\cref{sec:coefficient_mixed_audit}).  This empirical crossover accounts for
the intermediate fitted coefficient rate.  It does not contradict
\cref{prop:actual_ccop_local_defect}: the present diagnostic starts from an
accumulated numerical state, whereas that proposition is an exact-start
fixed-space result for a second-order matching predictor.
The exact-start action audit in \cref{tab:affine_stage_rates} applies to the
affine reference.  Here
\eqref{eq:test2_temporal_self_difference_rates} establishes the measured
fixed-cloud state self-convergence under nonaffine evolution.
The successive-difference curves are shown in
\cref{fig:supp_test2_temporal_self}.

\section{Discussion and conclusions}
\label{sec:discussion_conclusion}

The central result is a separation of two finite-step choices: terminal
configuration fixes the incompressibility target, whereas action configuration
fixes the pressure-correction range.  Their combination gives the scalar cross
response $D_TG_\Theta$ and its oblique projector.  Configuration assignment thus
extends the usual pressure-projection architecture without increasing the
pressure-space dimension or introducing another pressure stage.  The placement
parameter is geometric rather than merely temporal: it assigns the pressure
operator while leaving the terminal target fixed.

The distinction is measurable rather than merely algebraic.  With a fixed
terminal target, the affine calculations change the correction range and
pressure work by varying action placement.  The effect persists with a
mass-adjoint endpoint control, separating configuration-induced obliqueness
from the endpoint discretization baseline.  The exact-start stage test links
midpoint placement to cancellation of the leading configuration-induced
velocity defect.  Every placement satisfies the same exact discrete
work--energy identity, but the selected action determines which work is
represented and how accurately it approximates the physical pressure work.  In long-time
affine motion, the matched ISPH--CCOP comparison then transfers that one-step
mechanism to accumulated state, pressure, and energy consequences.

Nonaffine shear establishes a second, complementary distinction: consistency
of $G_h$, consistency of $D_hG_h$, and accuracy of the solved action are
different measurements.  Boundary propagation lowers the composite rate,
yet analytic-forcing pressure solves recover approximately second-order
actions for the three tested smooth pressures.  The decomposition in
\eqref{eq:solved_action_error_decomposition} identifies the inverse-transmitted
defect that connects these observations.  This provides a concrete diagnostic
for improving boundary actions without confusing algebraic closure with
spatial accuracy.

The analysis and calculations concern smooth two-dimensional inviscid flows.
Discrete closure is defined in the selected pressure-test space, while
Jacobian closure is measured independently.  Midpoint temporal accuracy
requires stage consistency, and the sampled response gaps provide
finite-dimensional stability evidence.  On finite nonaffine clouds, the
transported Piola rows leave a translation-equivariance defect that converges
as a spatial consistency error.
The matched ISPH comparison isolates a workflow change; comparisons with
higher-order time integrators require separate accuracy--cost measurements.
Within these conditions, CCOP supplies an explicit configuration choice whose
effect is resolved by pressure-work and deformation histories beyond the
terminal residual.  Equally accurate closure can therefore accompany distinct
finite-step action paths.

\appendix
\section{Exact-reference calculations}
\label{app:exact_references}

\subsection{Affine scalar reduction}

For $A=\diag(a,a^{-1})$,
$\dot A=\diag(\dot a,-\dot a/a^2)$ and
$\det\dot A=-\dot a^2/a^2$.  Since
$|A|_F^2=a^2+a^{-2}$,
\[
 \lambda=-\frac{2\det\dot A}{|A|_F^2}
 =\frac{2\dot a^2}{a^4+1}.
\]
The first diagonal component of $\ddot A=\lambda\cof A$ gives
$\ddot a=\lambda/a$.  Conservation of $|\dot A|_F^2/2=1$ selects the positive
branch in \eqref{eq:affine_scalar_reference}.

\subsection{Differential-vortex determinant and pressure}

Write
$F=\mathcal R(a)(I+(KX)\otimes\nabla a)$.  The matrix determinant lemma gives
$\det F=1+\nabla a\cdot KX=1$.  For
$u=\Omega(r)Kx$, the radial acceleration is
$(u\cdot\nabla)u=-\Omega(r)^2x$.  The momentum balance with $f=-x$ requires
$p_r=r(\Omega(r)^2-1)$.  Integrating from $r$ to 1 gives
\eqref{eq:differential_vortex_pressure}.

\section{Operator and reproducibility audits}

\subsection{Case 2: same-discretization ISPH implementation contract}
\label{sec:case2_implementation_contract}

The Case~2 archives contain 8037 particles, including 314 material boundary
labels and 7723 interior pressure coefficients.  They record $R=0.5$,
$\Delta x=0.01$, kernel length $0.013$, $\rho_0=1$, $A_0=0.4$, and
$\Omega=1.2$.  The strict drivers are
\path{card_piola_pressure_matrixfree_jax_v12_sph2_three_branch_strict.py}
and its \texttt{mls2} counterpart.  In the audited versions, both invoke
the same \texttt{build\_mls\_derivative\_operator} function for the material
derivatives used in $F,D$, and $G$.  The \texttt{SPH2} label identifies the
archive, not an independently verified conservative SPH pair.

For the unscaled basis
$[1,\xi,\eta,\xi^2/2,\xi\eta,\eta^2/2]$, the function uses cubic-kernel
weights within twice the kernel length, including the center.  A nearest-point
fallback supplies at least ten labels; a stencil with fewer than six
independent moments is rejected.  The derivative weights are extracted from
\begin{equation}
 (P_i^TW_iP_i+\lambda_iI)^{-1}P_i^TW_i,\qquad
 \lambda_i=\epsilon_{\rm mom}
 \left(\frac{\operatorname{tr}(P_i^TW_iP_i)}{6}+1\right).
 \label{eq:case2_moment_convention}
\end{equation}
The audited default is $\epsilon_{\rm mom}=10^{-12}$.  This is a local
moment regularization, distinct from a shift of the pressure matrix.
Consequently polynomial reproduction is subject to that perturbation and
conditioning, not an unconditional roundoff identity.  The reference
derivatives and boundary labels remain fixed throughout the run; the geometry
and Piola divergence in \eqref{eq:case2_physical_action} are assembled from
the current total material map.  They are not reconstructed from an unrelated
current-cloud fit.

The \texttt{star-star} branch solves the square frozen system
\eqref{eq:case2_same_discretization_isph}.  The \texttt{ccop-midpoint} branch
uses terminal divergence and midpoint physical-pressure action.  These are the
two workflows in the Case~2 main comparison.  The driver retains an optional
\texttt{terminal-terminal} development control, but it is not used to repeat
the target/action ablation already established in Case~1.  Selecting an action
position alone is not equivalent to selecting the star--star workflow, because
its target must also be $D_\star^q$.

The audited defaults use fp64, 30 unrolled geometry iterations per pressure
evaluation, automatic-differentiation JVPs, a relative nonlinear residual
target $10^{-9}$, a frozen-star linear target $10^{-10}$, and a rebuilt
kinematic RMS defect divided by $\Delta x$ below $10^{-10}$.  Accepted raw
Jacobians must exceed $10^{-6}$; the small safe determinant in trial evaluation
does not replace this raw accepted-state check.  These are distinct linear
and nonlinear acceptance rules, not equal iteration counts or equal cost.
In the stored midpoint runs, the maximum rebuilt relative residual is below
$4.1\times10^{-12}$ and the kinematic RMS defect is below $9.0\times10^{-17}$.
The NPZ files do not retain the historical source hash, moment parameter, or
complete solver configuration, so the current defaults cannot by themselves
certify the exact original command.  The source/data fingerprint audit records
this provenance boundary explicitly.

\subsection{Nonaffine coefficient self-difference crossover}
\label{sec:coefficient_mixed_audit}

For the three successive midpoint pressure differences at each fixed cloud,
the stored diagnostic fits the scalar RMS values to
$a_h\dt+b_h\dt^2$.  At $1/h=12,18,24$, respectively, the fitted coefficients
are
\[
 a_h=(1.1793,\,0.52165,\,0.29354)\times10^{-3},\qquad
 b_h=(0.044327,\,0.044287,\,0.044218).
\]
The effective rates are 1.234, 1.402, and 1.538; relative fit residuals are
below $1.9\times10^{-4}$.  The observed $h$-rate of $a_h$ is 2.007.
This three-point, two-parameter norm fit is consistent with a space--time
crossover.  It is not a vector expansion of the pressure or a proof that the
nonaffine multiplier automatically represents a time-averaged pressure.

\label{app:operator_audits}

The pressure-space and solver information omitted from the principal result
tables is retained here and in the machine-readable output.  Across the two
tested operator families, the square responses remain full rank and the
Riesz-scaled gaps in \cref{tab:supp_riesz_gap_summary} stay separated from
zero.  These are empirical bounds on the stated cloud and deformation
families, not a universal inf--sup claim.

\begin{table}[ht]
\centering
\caption{Minimum measured Riesz-scaled cross-response gaps.}
\label{tab:supp_riesz_gap_summary}
\begin{tabular}{lr}
\toprule
operator family & $\min\beta_{T,\Th,h}$ \\
\midrule
Test 1 affine pressure spaces & 0.7959 \\
Case 3 prescribed nonaffine maps & 0.4615 \\
\bottomrule
\end{tabular}
\end{table}

The affine disk
levels contain 423, 973, 1749, and 2751 interior pressure coefficients.  All
responses are full rank.  Their raw relative singular gaps decrease from
$3.85\times10^{-3}$ to $1.23\times10^{-3}$ and raw condition numbers increase
from 260 to 814; these unscaled quantities are reproducibility data, not the
stability norm used in the analysis.  The realized endpoint projector has
$\norm[M_U]{\Pi_T}=1.61863$ and
$e_{{\rm Gr},{\rm full}}=0.2484$ as defined in
\eqref{eq:full_space_green_compatibility}.  The latter is the relative
Frobenius mismatch of the full point-value pair, not an error estimate for
either differential operator.  On the resolved affine/quadratic subspaces,
the corresponding bilinear quantity
\eqref{eq:restricted_green_compatibility} is $2.26\times10^{-16}$.
Thus the endpoint projector is nonorthogonal on the full discrete space while
the polynomial solution subspace is compatible to roundoff.
The maximum $D_T$-anchored residual in the anchor sweep is
$4.18\times10^{-15}$.  Across the 21 frozen placements, the maximum relative
closure, idempotence, and action-annihilation defects are
$1.17\times10^{-15}$, $1.95\times10^{-15}$, and $9.41\times10^{-16}$,
respectively.  These values verify the algebra and direct solve; they are not
used as spatial-accuracy evidence.

In the differential-vortex prescribed-map study, the minimum Riesz-scaled
cross gap over all configurations is 0.4615.  The exact-increment transport
depth audit is summarized in \cref{fig:supp_transport}.

\begin{table}[ht]
\centering
\caption{Prescribed-map midpoint audit at the stronger shear $\chi=1.5$.
Here $\max\kappa_2(F_{\rm ex})=10.91$ and the exact relative-deformation
margin is $0.25$ on every level.}
\label{tab:supp_strong_shear}
\begin{tabular}{rrrrrr}
\toprule
$R/h$ & $e_F$ & $e_G$ & $e_D^q$ & $e_{DG}^q$ & $\beta_{T,1/2,h}$ \\
\midrule
12 & $5.374\mathrm e{-2}$ & $8.508\mathrm e{-3}$ & $1.610\mathrm e{-2}$ & $8.167\mathrm e{-2}$ & 0.653 \\
18 & $2.363\mathrm e{-2}$ & $3.816\mathrm e{-3}$ & $7.626\mathrm e{-3}$ & $5.129\mathrm e{-2}$ & 0.973 \\
24 & $1.325\mathrm e{-2}$ & $2.153\mathrm e{-3}$ & $4.504\mathrm e{-3}$ & $3.615\mathrm e{-2}$ & 1.260 \\
30 & $8.464\mathrm e{-3}$ & $1.377\mathrm e{-3}$ & $2.985\mathrm e{-3}$ & $2.727\mathrm e{-2}$ & 1.246 \\
\bottomrule
\end{tabular}
\end{table}

The constant-field check in \cref{sec:constant_field_diagnostics} was applied
to the same prescribed nonaffine map at $\chi=1$ and to the midpoint action
path.  Table~\ref{tab:supp_constant_translation_audit} reports the larger of
the two Cartesian constant-vector defects.  The second column is a terminal
constraint-space RMS; the third is the mass-weighted velocity defect of the
frozen projector under addition of a unit translation.  No corrective
projection is applied before either measurement.

\begin{table}[ht]
\centering
\caption{Constant-vector and frozen translation-equivariance audit for the
prescribed nonaffine map at $\chi=1$, $\Th=1/2$.}
\label{tab:supp_constant_translation_audit}
\begin{tabular}{rrr}
\toprule
$1/h$ & $\max_k\norm[M_Z,{\rm rms}]{D_T^qe_k}$
      & $\max_k\norm[M_U,{\rm rms}]{G_\Th(D_T^qG_\Th)^{-1}D_T^qe_k}$ \\
\midrule
12 & $1.181\times10^{-2}$ & $1.195\times10^{-3}$ \\
18 & $7.054\times10^{-3}$ & $4.915\times10^{-4}$ \\
24 & $4.900\times10^{-3}$ & $2.628\times10^{-4}$ \\
30 & $3.504\times10^{-3}$ & $1.572\times10^{-4}$ \\
\midrule
fitted rate & 1.315 & 2.208 \\
\bottomrule
\end{tabular}
\end{table}

The pressure solve and projector identities remain algebraically exact for
their represented rows, but the first column confirms that Piola matrix
transport alone does not create an exact discrete geometric-conservation law.
The induced frozen translation error nevertheless decreases approximately
quadratically over this cloud family.  This is an operator-level translation
check.  Because the differential vortex contains the origin-fixed force
$f=-x$, it is not itself a Galilean-invariant trajectory benchmark.

\begin{figure}[ht]
\centering
\safegraphic[0.72\linewidth]{\figroot/test2_differential_vortex/figures/part2b_transport_depth.png}
\caption{Exact-increment transport audit.  Recursive/direct operator defects
remain at roundoff through 64 subintervals, whereas direct spatial
reconstruction sets the visible error.}
\label{fig:supp_transport}
\end{figure}

\begin{figure}[ht]
\centering
\safegraphic[0.88\linewidth]{\figroot/test2_differential_vortex/figures/part2c_temporal_self_convergence.png}
\caption{Fixed-cloud temporal self differences for the complete nonaffine
evolution.  Midpoint state and geometry differences follow the second-order
path, whereas terminal differences follow the first-order path; the exact
error itself is spatial-floor limited on this cloud.}
\label{fig:supp_test2_temporal_self}
\end{figure}

The deterministic seed is 20260728 for Cases~1 and~3.  The principal drivers
and auxiliary audit scripts are
\begin{itemize}[leftmargin=2em]
\item \path{test1_affine_free_boundary.py},
\item \path{card_piola_pressure_matrixfree_jax_v12_sph2_three_branch_strict.py}
and its \texttt{mls2} counterpart for Case~2,
\item \path{test2_differential_vortex.py},
\item \path{constant_translation_audit.py}.
\end{itemize}
The drivers generate complete JSON and per-iteration histories.  Compact CSV
tables, configuration records, reports, and figures are archived below
\path{results/test1_affine_free_boundary},
\path{results/test2_affine_long_time},
and \path{results/test2_differential_vortex}.  The unified regression command is
\begin{verbatim}
python -m pytest test_test1_affine_free_boundary.py \
  test_test2_differential_vortex.py -q
\end{verbatim}

\begin{thebibliography}{99}


\bibitem{MarsdenWest2001}
J.~E. Marsden and M.~West,
Discrete mechanics and variational integrators,
\emph{Acta Numer.}, 10 (2001), pp.~357--514.

\bibitem{Andersen1983}
H.~C. Andersen,
RATTLE: A ``velocity'' version of the SHAKE algorithm for molecular dynamics
calculations,
\emph{J. Comput. Phys.}, 52 (1983), pp.~24--34,
\url{https://doi.org/10.1016/0021-9991(83)90014-1}.

\bibitem{Betsch2005}
P.~Betsch,
The discrete null space method for the energy consistent integration of
constrained mechanical systems: Part I: Holonomic constraints,
\emph{Comput. Methods Appl. Mech. Engrg.}, 194 (2005), pp.~5159--5190,
\url{https://doi.org/10.1016/j.cma.2005.01.004}.

\bibitem{SuzukiKoshizukaOka2007}
Y.~Suzuki, S.~Koshizuka, and Y.~Oka,
Hamiltonian moving-particle semi-implicit (HMPS) method for incompressible
fluid flows,
\emph{Comput. Methods Appl. Mech. Engrg.}, 196 (2007), pp.~2876--2894,
\url{https://doi.org/10.1016/j.cma.2006.12.006}.


\bibitem{ElleroSerranoEspanol2007}
M.~Ellero, M.~Serrano, and P.~Espa\~nol,
Incompressible smoothed particle hydrodynamics,
\emph{J. Comput. Phys.}, 226 (2007), pp.~1731--1752,
\url{https://doi.org/10.1016/j.jcp.2007.06.019}.

\bibitem{NairTomar2015}
P.~Nair and G.~Tomar,
Volume conservation issues in incompressible smoothed particle hydrodynamics,
\emph{J. Comput. Phys.}, 297 (2015), pp.~689--699,
\url{https://doi.org/10.1016/j.jcp.2015.05.042}.

\bibitem{BasicEtAl2022}
J.~Ba\v{s}i\'c, N.~Degiuli, B.~Blagojevi\'c, and D.~Ban,
Lagrangian differencing dynamics for incompressible flows,
\emph{J. Comput. Phys.}, 462 (2022), 111198.

\bibitem{GotohKhayyerGotoh2024}
T.~Gotoh, A.~Khayyer, and H.~Gotoh,
Enhanced schemes for resolution of the continuity equation in
projection-based SPH,
\emph{Engrg. Anal. Bound. Elem.}, 166 (2024), 105848.

\bibitem{Matsunaga2025}
T.~Matsunaga,
High-order time-marching schemes for incompressible flow in particle methods,
\emph{Comput. Methods Appl. Mech. Engrg.}, 447 (2025), 118395.

\bibitem{MatsunagaKoshizuka2022}
T.~Matsunaga and S.~Koshizuka,
Stabilized LSMPS method for complex free-surface flow simulation,
\emph{Comput. Methods Appl. Mech. Engrg.}, 389 (2022), 114416.

\bibitem{SuchdeEtAl2025}
P.~Suchde, C.~Leith\"auser, J.~Kuhnert, and S.~P.~A. Bordas,
Volume and mass conservation in Lagrangian meshfree methods,
\emph{Int. J. Numer. Methods Engrg.}, 126 (2025), e7657.

\bibitem{ZhangEtAl2024}
N.~Zhang, S.~Yan, Q.~Ma, A.~Khayyer, X.~Guo, and X.~Zheng,
A consistent second order ISPH for free surface flow,
\emph{Comput. Fluids}, 274 (2024), 106224.

\bibitem{AlmgrenBellCrutchfield2000}
A.~S. Almgren, J.~B. Bell, and W.~Y. Crutchfield,
Approximate projection methods: Part I. Inviscid analysis,
\emph{SIAM J. Sci. Comput.}, 22 (2000), pp.~1139--1159.

\bibitem{AlmgrenBellSzymczak1996}
A.~S. Almgren, J.~B. Bell, and W.~G. Szymczak,
A numerical method for the incompressible Navier--Stokes equations based on an
approximate projection,
\emph{SIAM J. Sci. Comput.}, 17 (1996), pp.~358--369.

\bibitem{Chorin1968}
A.~J. Chorin,
Numerical solution of the Navier--Stokes equations,
\emph{Math. Comp.}, 22 (1968), pp.~745--762.

\bibitem{CumminsRudman1999}
S.~J. Cummins and M.~Rudman,
An SPH projection method,
\emph{J. Comput. Phys.}, 152 (1999), pp.~584--607.

\bibitem{DiLiTangZhang2005}
Y.~Di, R.~Li, T.~Tang, and P.~Zhang,
Moving mesh finite element methods for the incompressible Navier--Stokes
equations,
\emph{SIAM J. Sci. Comput.}, 26 (2005), pp.~1036--1056.

\bibitem{DoneaHuerta2004}
J.~Donea, A.~Huerta, J.-P. Ponthot, and A.~Rodriguez-Ferran,
Arbitrary Lagrangian--Eulerian methods,
in \emph{Encyclopedia of Computational Mechanics}, Wiley, Chichester, 2004.

\bibitem{LesoinneFarhat1996}
M.~Lesoinne and C.~Farhat,
Geometric conservation laws for flow problems with moving boundaries and
deformable meshes, and their impact on aeroelastic computations,
\emph{Comput. Methods Appl. Mech. Engrg.}, 134 (1996), pp.~71--90,
\url{https://doi.org/10.1016/0045-7825(96)01028-6}.

\bibitem{GuillardFarhat2000}
H.~Guillard and C.~Farhat,
On the significance of the geometric conservation law for flow computations
on moving meshes,
\emph{Comput. Methods Appl. Mech. Engrg.}, 190 (2000), pp.~1467--1482,
\url{https://doi.org/10.1016/S0045-7825(00)00173-0}.

\bibitem{FormaggiaNobile2004}
L.~Formaggia and F.~Nobile,
Stability analysis of second-order time accurate schemes for ALE--FEM,
\emph{Comput. Methods Appl. Mech. Engrg.}, 193 (2004), pp.~4097--4116,
\url{https://doi.org/10.1016/j.cma.2003.09.028}.

\bibitem{PavlovEtAl2011}
D.~Pavlov, P.~Mullen, Y.~Tong, E.~Kanso, J.~E. Marsden, and M.~Desbrun,
Structure-preserving discretization of incompressible fluids,
\emph{Physica D}, 240 (2011), pp.~443--458,
\url{https://doi.org/10.1016/j.physd.2010.10.012}.

\bibitem{GawlikEtAl2011}
E.~S. Gawlik, P.~Mullen, D.~Pavlov, J.~E. Marsden, and M.~Desbrun,
Geometric, variational discretization of continuum theories,
\emph{Physica D}, 240 (2011), pp.~1724--1760,
\url{https://doi.org/10.1016/j.physd.2011.07.011}.

\bibitem{GuermondMinevShen2006}
J.-L. Guermond, P.~Minev, and J.~Shen,
An overview of projection methods for incompressible flows,
\emph{Comput. Methods Appl. Mech. Engrg.}, 195 (2006), pp.~6011--6045.

\bibitem{KhayyerGotoh2011}
A.~Khayyer and H.~Gotoh,
Enhancement of stability and accuracy of the moving particle semi-implicit
method,
\emph{J. Comput. Phys.}, 230 (2011), pp.~3093--3118.

\bibitem{KoshizukaOka1996}
S.~Koshizuka and Y.~Oka,
Moving-particle semi-implicit method for fragmentation of incompressible
fluid,
\emph{Nucl. Sci. Eng.}, 123 (1996), pp.~421--434.

\bibitem{MirzaeiSchabackDehghan2012}
D.~Mirzaei, R.~Schaback, and M.~Dehghan,
On generalized moving least squares and diffuse derivatives,
\emph{IMA J. Numer. Anal.}, 32 (2012), pp.~983--1000.

\bibitem{RobertsShkollerSideris2020}
J.~Roberts, S.~Shkoller, and T.~C. Sideris,
Affine motion of 2D incompressible fluids surrounded by vacuum and flows in
$\mathrm{SL}(2,\R)$,
\emph{Comm. Math. Phys.}, 375 (2020), pp.~1003--1040.

\bibitem{Temam1969}
R.~Temam,
Sur l'approximation de la solution des equations de Navier--Stokes par la
methode des pas fractionnaires II,
\emph{Arch. Ration. Mech. Anal.}, 33 (1969), pp.~377--385.

\bibitem{TraskMaxeyHu2018}
N.~Trask, M.~Maxey, and X.~Hu,
A compatible high-order meshless method for the Stokes equations with
applications to suspension flows,
\emph{J. Comput. Phys.}, 355 (2018), pp.~310--326.

\bibitem{TraskPeregoBochev2017}
N.~Trask, M.~Perego, and P.~Bochev,
A high-order staggered meshless method for elliptic problems,
\emph{SIAM J. Sci. Comput.}, 39 (2017), pp.~A479--A502.

\bibitem{XuStansbyLaurence2009}
R.~Xu, P.~Stansby, and D.~Laurence,
Accuracy and stability in incompressible SPH (ISPH) based on the projection
method and a new approach,
\emph{J. Comput. Phys.}, 228 (2009), pp.~6703--6725.


\bibitem{Wang2019IPST}
P.-P.~Wang, Z.-F.~Meng, A.-M.~Zhang, F.-R.~Ming, and P.-N.~Sun,
Improved particle shifting technology and optimized free-surface detection
method for free-surface flows in smoothed particle hydrodynamics,
\emph{Comput. Methods Appl. Mech. Engrg.}, 357 (2019), 112580,
\url{https://doi.org/10.1016/j.cma.2019.112580}.

\bibitem{ZhangHighOrderSPH}
B.~Zhang, N.~A. Adams, and X.~Y. Hu,
Towards high-order consistency and convergence of conservative SPH
approximations,
\emph{Comput. Methods Appl. Mech. Engrg.}, 433 (2025), 117484,
\url{https://doi.org/10.1016/j.cma.2024.117484}.

\end{thebibliography}

\end{document}
